\chapter{Equations over Finite Fields}

The central question of this chapter is: given a polynomial $f\in\bbF_q[X_1,\dots,X_n]$, how many solutions does $f(X_1,\dots,X_n)=0$ have in $\bbF_q^n$? We write $\mcV(f)$ for the solution set and $Z(f)=|\mcV(f)|$ for its size; for a system $f_1=\cdots=f_k=0$ we write $\mcV(f_1,\dots,f_k)=\mcV(f_1)\cap\cdots\cap\mcV(f_k)$ and $Z(f_1,\dots,f_k)$ for the common solution count. Variants $\ov{\mcV}$ and $\ov{Z}$ track the non-trivial (nonzero) solutions. The notation mirrors that of varieties of algebraic geometry: $\mcV(f)$ is the affine variety cut out by $f$.

The chapter develops two complementary strategies. The first is \emph{direct combinatorics}: for special polynomial families (univariate polynomials, systems of low total degree, quadratic forms, diagonal equations) one can write exact formulas or prove sharp divisibility constraints on $Z(f)$. The second is \emph{character sum methods}: the indicator of whether $\alpha$ is a solution can be encoded as an additive character sum, turning the problem of counting solutions into the problem of bounding $\sum_\alpha\chi(f(\alpha))$. These two strategies reinforce each other throughout the chapter.

\section{Some Elementary Results}
In this section we give some elementary results on the number of solutions. We present some classical theorems such as those of K\"{o}nig-Rados, Chevalley, and Warning. We also establish elementary upper bounds for the number of solutions and results on the expected order of magnitude.
\subsection{Univariate Polynomials}
We start of with univariate polynomials. Let $f\in\bbF_q[X]$ is an univariate polynomial Naturally the number of solutions in $\bbF_q$ for the equation $f(X)=0$ is basically the degree of $\gcd(f(X),X^q-X)$. Since checking if $0$ is a solution or not is very easy to check we'll only consider non-zero solutions. So the number of non-zero solutions of $f(X)=0$ is basically the degree of $\gcd(f(X),X^{q-1}-1)$. Therefore we may always assume without loss of generality that $\deg(f)\leq q-1$. Suppose $$f(X)=a_{q-1}\cdot X^{q-1}+a_{q-2}\cdot X^{q-2}+\cdots+a_1\cdot X+a_0$$where $a_i\in\bbF_q$ for all $0\leq i\leq q-1$. Since for all $\alpha\in\bbF_q^*$ we have $\alpha^{q-1}=1$ the variable $X^{q-1}$ will always evaluate to be $1$ for all non-zero values of $X$. Therefore the number of solutions of $f$ is same as of the equation $a_{q-2}\cdot X^{q-2}+\cdots+a_1\cdot X+(a_0+a_{q-1})=0$. Therefore we can always assume $\deg(f)\leq q-2$. So suppose we have a polynomial $f\in\bbF_q[X]$ such that $$f(X)=a_{q-2}\cdot X^{q-2}+\cdots+a_1\cdot X+a_0$$where $a_i\in\bbF_q$ for all $0\leq i\leq q-2$. 

Now we can also analyze the number of solutions of $f$ using matrix methods. From $f$ we can define the following matrix $$M_f=\mat{a_0 & a_1 & \cdots & a_{q-2}\\ a_1 & a_2 & \cdots & a_0\\ \vdots & \vdots & \ddots & \vdots\\ a_{q-2} & a_0 & \cdots & a_{q-3}}$$This matrix is a \emph{left circulant matrix} i.e. each row is obtained from the preceding row by a cyclic shift of the entries to the left. The number of non-zero solutions of $f(X)=0$ and the rank of the matrix $M_f$ are very related. 
\begin{Theorem}{K\"{o}nig-Rados Theorem}{thm:konig-rados}
  Let $f(X)=\sum_{i=0}^{q-2}a_i\cdot X^i\in\bbF_q[X]$ be an univariate polynomial. Then $$\ov{Z}(f)=q-1-\Rank(M_f)$$  
\end{Theorem}
\begin{proof}
  Since $\deg(f)\leq q-2$ we have $\ov{Z}(f)\leq q-2$. $\ov{Z}(f)=q-1$ rises only when $f$ is a zero polynomial i.e. $M_f$ is a zero matrix and in that case $\Rank(M_f)=0$ which satisfies the theorem. Let $\gm$ be the primitive element of $\bbF_q$. Then consider the \emph{Vandermonde Matrix}, $V(\gm,\gm^2,\dots, \gm^{q-1})$: $$V=\mat{1&1&\cdots&1\\ \gm & \gm^2 &\cdots & \gm^{q-1}\\ \vdots & \vdots & \ddots & \vdots\\ \gm^{q-2} & \gm^{2 (q-2)} & \cdots &\gm^{(q-1)(q-2)}}$$Then using $\gm^{i\cdot (q-1)}=1$ for all $i\in[q-1]$ we obtain $$M_f\cdot V=\mat{f(\gm) & f(\gm^2) & \cdots & f(\gm^{q-1})\\ \gm^{-1}f(\gm) & \gm^{-2} f(\gm^2) & \cdots & \gm^{-(q-1)}f(\gm^{q-1})\\ \vdots & \vdots & \ddots & \vdots \\ \gm^{-(q-2)}f(\gm) & \gm^{-2(q-2)}f(\gm^2) & \cdots & \gm^{-(q-1)(q-2)}f(\gm^{q-1})}$$Now for any $i\in[q-1]$ we have $f(\gm^i)=0$ if and only if  the $i^{th}$ row becomes zero. Therefore the $\Rank(M_f\cdot V)=q-1-\ov{Z}(f)$. Since $V$ is the Vandermonde matrix $V$ and $M_f\cdot V$ have the same rank. So the theorem follows. 
\end{proof}

In the next subsections we will give simple upper and lower bounds for solutions multivariate polynomial equations or system of polynomial equations. 
\subsection{Chevalley--Warning Theorem}
We begin with two elementary lemmas about power sums that will serve as the key tools.
\begin{lemma}{}{lem:power-sum}
  Let $k$ be a non-negative integer. Suppose for $k=0$, $0^0=1$. Then $$\sum_{\alpha\in\bbF_q}\alpha^k=\begin{cases}
    0 & \text{if $k=0$ or $q-1\nmid k$}\\ 
    -1 & \text{if $q-1\mid k$}
  \end{cases}$$
\end{lemma}
\begin{proof}
  For $k=0$ the sum becomes $q\equiv 0 \bmod q$. So suppose $k>0$. Let $\gm$ is the primitive element of $\bbF_q$. Then we have $$\sum_{\alpha\in\bbF_q}\alpha^k=\sum_{\alpha\in\bbF_q^*}\alpha^k=\sum_{j=0}^{q-2}\gm^{j\cdot k}=\sum_{j=0}^{q-2}(\gm^k)^j$$Now if $q-1\mid k$ then $\gm^k=1$ so the sum evaluates to be $-1$. If $q-1\nmid k$ then $\sum_{j=0}^{q-2}(\gm^k)^j=\frac{\gm^{k(q-1)}-1}{\gm^k-1}=0$. Hence we have the lemma. 
\end{proof}
So using this lemma we can prove if we take the sum over all evaluations of a polynomial then the sum evaluates to be zero.
\begin{lemma}{}{lem:poly-evaluate-allsum}
  Let $f\in\bbF_q[X_1,\dots,X_n]$ with $\deg(f)<n(q-1)$. Then $$\sum_{\alpha_1,\dots,\alpha_n\in\bbF_q}f(\alpha_1,\dots, \alpha_k)=0$$
\end{lemma}
\begin{proof}
  It suffices to prove this property for every monomial of $f$. So suppose $X_1^{k_1}\cdots X_n^{k_n}$ be a monomial of $f$ where $k_1+\cdots+k_n<n(q-1)$ and $k_i\in\bbZ_0$ for all $i\in[n]$. Since $k_1+\cdots+k_n<n(q-1)$ then by Pigeon Hole Principle there exists $j\in[n]$ such that $k_j< q-1$. Then
  $$\sum_{\alpha_1,\dots,\alpha_n\in\bbF_q}\alpha_1^{k_1}\cdots \alpha_n^{k_n}=\prod_{i=1}^n \lt(\sum_{\alpha_i\in\bbF_q}\alpha^{k_i}\rt)$$ Then for $j$, by \lemref{lem:power-sum} $\sum_{\alpha_j\in\bbF_q}\alpha^{k_j}=0$. So we have the lemma.
\end{proof}

The two lemmas above are the principal tools behind the Chevalley--Warning theorem. It asserts that whenever the total degree of a system is small relative to the number of variables, the number of common zeros is divisible by $p$ --- forcing, in particular, that a system with one solution must have a second.

\begin{Theorem}{Chevalley--Warning Theorem}{thm:chevalley-warning}
  \index{Chevalley-Warning theorem}%
  Let $f\in\bbF_q[X_1,\dots,X_n]$ with $\deg(f)<n$. Then the number of solutions of the equation $f(X_1,\dots, X_n)=0$ in $\bbF_q^n$ is divisible by the characteristic $p$ of $\bbF_q$. 

  Moreover, if $f(0,\dots, 0)=0$ then there exists a nontrivial solution in $\mcV(f)$ i.e. $\ov{\mcV}(f)\neq \emptyset$ or $\ov{Z}(f)\geq 1$. 
\end{Theorem}
\begin{proof}
  We first construct a polynomial and then will use \PolyEvallem. So take the polynomial $F(X_1,\dots, X_n)=1-f^{q-1}(X_1,\dots, X_n)$. By construction $F$ always takes vale $0$ or $1$. Therefore for any $\alpha_1,\dots, \alpha_n\in\bbF_q$ we have $$F(\alpha_1,\dots,\alpha_n)=1\iff f(\alpha_1,\dots, \alpha_n)=0$$So $$\sum_{\alpha_1,\dots,\alpha_n\in\bbF_q}F(\alpha_1,\dots, \alpha_k)=Z(f)\bmod p$$Now by \PolyEvallem we have $$\sum_{\alpha_1,\dots,\alpha_n\in\bbF_q}F(\alpha_1,\dots, \alpha_k)=0$$Therefore $p\mid Z(f)$. So we have the first part. 

  Now since $f(0,\dots, 0)=0$ we have $Z(f)\geq 1$. Since $p$ is a prime $p\geq 2$. So by first part $p\mid Z(f)$. Hence $Z(f)\geq 2$. So we have the second part too. 
\end{proof}

Now the bound of $\deg(f)<n$ in \CWthm is the best possible bound for the theorem to work. We can construct a polynomial in $\bbF_q[X_1,\dots, X_n]$ of degree $n$ for which the theorem doesn't hold. Let $\{\alpha_1,\dots, \alpha_n\}$ is the basis of $E=\bbF_{q^n}$ over $\bbF_q$. Consider the following polynomial $$f(X_1,\dots, X_n)=\prod_{j=0}^{n-1}\lt(\alpha_1^{q^i}X_1+\cdots +\alpha_n^{q^i}X_n\rt)=\prod_{j=0}^{n-1}\sum_{j=1}^n \alpha_j^{q^i}X_i$$Since for all $i\in\{0,1,\dots, n-1\}$, $\alpha_j^{q^i}$ are conjugates of $\alpha_j$ for all $j\in[n]$ with respect to $\bbF_q$, $f$ is a polynomial over $\bbF_q$. Suppose for any $(\gm_1,\dots, \gm_n)\in\bbF_q^n$ consider the element $\alpha_1\cdot \gm_1+\cdots+\alpha_n\cdot \gm_n\in E$. Then we evaluate $f$ on $(\gm_1,\dots, \gm_n)$ and we get $$f(\gm_1,\dots,\gm_n)=\prod_{i=0}^{n-1}\sum_{j=1}^n \alpha_j^{q^i}\cdot \gm_j=\prod_{i=0}^{n-1}\lt(\sum_{j=1}^n \alpha_j\cdot \gm_j\rt)^{q^i}=N_{E/\bbF_q}(\alpha_1\cdot \gm_1+\cdots+\alpha_n\cdot \gm_n)$$Hence $f(\gm_1,\dots,\gm)=0$ if and only if $N_{E/\bbF_q}(\alpha_1\cdot \gm_1+\cdots+\alpha_n\cdot \gm_n)=0$ which holds only for $\alpha_1\cdot \gm_1+\cdots+\alpha_n\cdot \gm_n=0$. Since  $\{\alpha_1,\dots, \alpha_n\}$ forms a basis of $E$ over $\bbF_q$ we have $f(\gm_1,\dots,\gm)=0$ if and only if $\gm_j=0$ for all $j\in[n]$. Hence $Z(f)=1$. The \CWthm clearly doesn't hold in this case. 

Now we can extend the \CWthm to a system of polynomial equations $f_1=0,\dots,f_k=0$ with $f_1,\dots,f_k\in\bbF_q[X_1,\dots,X_n]$. In this case we are interested in the number of common solutions. 
\begin{Theorem}{}{}
  Let $f_1,\dots, f_k\in\bbF_q[X_1,\dots, X_n]$ with $\deg(f_1\cdots f_k)=\deg(f_1)+\cdots+\deg(f_k)<n$. Then $p\mid Z(f_1,\dots, f_k)$ where $p=\Char(\bbF_q)$. 

  Moreover, if $(0,\dots, 0)\in \mcV(f_1,\dots, f_k)$ then $\ov{\mcV}(f_1,\dots, f_k)\neq\emptyset$ i.e. $\ovZ(f_1,\dots, f_k)\geq 1$. 
\end{Theorem}
\begin{proof}
  Again we construct a polynomial. Consider the polynomial $F\in\bbF_q[X_1,\dots, X_n]$ such that $F=(1-f_1^{q_1})\cdots (1-f_k^{q-1})$. Therefore for any $\alpha_1,\dots,\alpha_n\in\bbF_q$ we have $F(\alpha_1,\dots,\alpha_n)=1$ if and only if $f_i(\alpha_1,\dots, \alpha_n)=0$ for all $i\in[k]$. Now $\deg(F)\leq \deg(f_1)+\cdots+\deg(f_k)<n(q-1)$. Now by applying \CWthm for $F$ we get the theorem. 
\end{proof}

\subsection{Lower Bounds on the Number of Solutions}

In the next theorem we will show if we have a system of polynomial equations $f_1=0,\dots,f_k=0$ with $f_1,\dots,f_k\in\bbF_q[X_1,\dots,X_n]$ then for any affine subspace $W$ all its parallel shifts have same number of solutions modulo the characteristic of $\bbF_q$. If $W_1$ and $W_2$ are two affine subspaces of $\bbF_q^n$ of same dimension then we call they are parallel affine subspaces if they are obtained by translation from the same linear subspace. 
\begin{lemma}{}{lem:solutions-affine-subspace-parallel-shift}
  Let $f_1,\dots, f_k\in\bbF_q[X_1,\dots, X_n]$.  If $W_1$ and $W_2$ are two parallel affine subspaces of $\bbF_q^n$ of dimension $d$ where $d=\deg(f_1)+\cdots+\deg(f_k)$ then $$|W_1\cap \mcV(f_1,\dots,f_k)|\equiv |W_2\cap \mcV(f_1,\dots, f_k)|\bmod p$$where $p=\Char(\bbF_q)$. 
\end{lemma}
\begin{proof}
  Now if $W_1=W_2$, we get $|W_1\cap \mcV(f_1,\dots,f_k)|=|W_2\cap \mcV(f_1,\dots, f_k)|$. So we get the theorem. So we can assume $W_1\neq W_2$. Now by change of coordinates we can take \begin{align*}
    & W_1 = \{(\alpha_1,\dots, \alpha_n)\in\bbF_q^n\colon \alpha_i=0,\ \forall\ i\in[n]\}\\ 
    & W_2 = \{(\alpha_1,\dots, \alpha_n)\in\bbF_q^n\colon \alpha_1=1\text{ and }\alpha_i=0,\ \forall\ 2\leq i\leq n\}\\ 
  \end{align*}
  Again we construct a polynomial $G\in\bbF_q[X_1,\dots, X_n]$ where $$G(X_1,\dots, X_n)=(-1)^{n-d}(X_1^{q-2}+\cdots X_1+1)(X_2^{q-1}-1)\cdots (X_n^{q-1}-1)$$So $\deg(G)=(n-d)(q-1)-1$. 
  \begin{observation}
    With the above construction of $G$ for any $(\alpha_1,\dots,\alpha_n)\in\bbF_q^n$ \begin{enumerate}[label=(\roman*)]
      \item if $(\alpha_1,\dots,\alpha_n)\in W_1$ then $G(\alpha_1,\dots, \alpha_n)=-1$
      \item if $(\alpha_1,\dots,\alpha_n)\in W_2$ then $G(\alpha_1,\dots, \alpha_n)=1$
      \item if $(\alpha_1,\dots,\alpha_n)\notin W_1\cup W_2$ then $G(\alpha_1,\dots, \alpha_n)=0$
    \end{enumerate}
  \end{observation}
  Then we construct the following polynomial $F\in\bbF_q[X_1,\dots, X_n]$ such that $$F=(1-f_1^{q-1})\cdots (1-f_k^{q-1})G$$So now we have $\deg(F)\leq d(q-1)+(n-d)(q-1)-1=n(q-1)-1$. Therefore for any $(\alpha_1,\dots,\alpha_n)\in\bbF_q^n$, if $(\alpha_1,\dots,\alpha_n)\in W_1\cap \mcV(f_1,\dots, f_k)$ then $F(\alpha_1,\dots, \alpha_n)=1$. If $(\alpha_1,\dots,\alpha_n)\in W_2\cap \mcV(f_1,\dots, f_k)$ then $F(\alpha_1,\dots, \alpha_n)=-1$ and $0$ elsewhere. So $$\sum_{\alpha_1,\dots,\alpha_n\in\bbF_q}F(\alpha_1,\dots,\alpha_n)=|W_1\cap \mcV(f_1,\dots,f_k)|- |W_2\cap \mcV(f_1,\dots, f_k)|$$Now by using \CWthm on $F$ we have the lemma. 
\end{proof}

The \AffSubspaceShiftlem above can be used to obtain a lower bound on $Z(f_1,\dots,f_k)$: the \CWthm guarantees that $p\mid Z(f_1,\dots,f_k)$ whenever $d<n$, but by counting how many parallel affine hyperplanes each must contain a solution we can say much more.
\begin{Theorem}{}{}
  Let $f_1,\dots, f_k\in\bbF_q[X_1,\dots, X_n]$ with $d=\deg(f_1)+\cdots+\deg(f_k)<n$. If $Z(f_1,\dots, f_k)\geq 1$ then $Z(f_1,\dots, f_k)\geq q^{n-d}$
\end{Theorem}
\begin{proof}
  We will prove this by analyzing two cases. 
  \paragraph*{Case I:} Let there exists an affine subspace $W_1$ of $\bbF_q^n$ with dimension $d$ such that $|W_1\cap \mcV(f_1,\dots, f_k)|\not\equiv 0\bmod p$. Then by \AffSubspaceShiftlem for all parallel affine spaces $W_2$ of $W_1$ we have $|W_2\cap \mcV(f_1,\dots,f_k)|\not\equiv 0\bmod p$. In particular $|W_2\cap \mcV(f_1,\dots, f_k)|\geq 1$. Since $W_1$ is of dimension $d$ there are $q^{n-d}$ distinct affine subspaces parallel to $W_1$. Therefore 
  $$\mcV(f_1,\dots, f_k)=\bigsqcup_{\substack{W\colon\text{affine subspace}\\ \text{parallel to $W_1$}}}\Big(W\cap \mcV(f_1,\dots,f_k)\Big)$$
  For each such $W$ we have $|W\cap\mcV(f_1,\dots,f_k)|\geq 1$. Therefore 
  $$Z(f_1,\dots, f_k)=\sum_{\substack{W\colon\text{affine subspace}\\ \text{parallel to $W_1$}}}|W\cap\mcV(f_1,\dots,f_k)|\geq q^{n-d}$$
  Therefore we have the theorem in this case. 
  \paragraph*{Case II:} We have $|W\cap\mcV(f_1,\dots,f_k)|\equiv 0\bmod p$. for all affine subspaces $W$ of $\bbF_q^n$ of dimension $d$. Since $Z(f_1,\dots,f_k)\geq 1$ there exists $ k\in[d]$ such that for any affine subspace $W'$ of dimension $k$, $|W'\cap \mcV(f_1,\dots,f_k)|\equiv 0\bmod p$ but there is an affine subspace $W''$ of dimension $(k-1)$ such that $|W''\cap \mcV(f_1,\dots, f_k)|\not\equiv 0\bmod p$. Now consider all affine subspaces $W'$ of dimension $k$ containing $W''$. There are exactly $\frac{q^{n-(k-1)}-1}{q-1}$ many such affine subspaces. For each such affine subspace $W'$ we have $$|(W'\setminus W'')\cap \mcV(f_1,\dots,f_k)|=|W'\cap \mcV(f_1,\dots,f_k)|-|W''\cap \mcV(f_1,\dots,f_k)|\not\equiv0\bmod p$$So $|(W'\setminus W'')\cap \mcV(f_1,\dots,f_k)|\geq 1$. Therefore we have $$\mcV(f_1,\dots, f_k)=\Big(W''\cap \mcV(f_1,\dots,f_k)\Big)\sqcup\bigsqcup_{\substack{W'\colon \text{affine subspace}\\ \text{of dimension $k$}\\ \text{contains $W''$}}}\Big((W'\setminus W'')\cap \mcV(f_1,\dots,f_k)\Big)$$Hence we have $$Z(f_1,\dots, f_k)=|W''\cap \mcV(f_1,\dots,f_k)|+\sum_{\substack{W'\colon \text{affine subspace}\\ \text{of dimension $k$}\\ \text{contains $W''$}}}|(W'\setminus W'')\cap \mcV(f_1,\dots,f_k)|\geq 1+\frac{q^{n-(k-1)}-1}{q-1}>q^{n-d}$$So we have the theorem. 
\end{proof}

The lower bound on the number of solutions above is best possible even for $k=1$ i.e. for any positive integers $d$ and $n$ with $d<n$ there is a polynomial $f\in\bbF_q[X_1,\dots, X_n]$ of degree $d$ such that $f(X_1,\dots, X_n)=0$ has exactly $q^{n-d}$ solutions in $\bbF_q^n$. Consider the polynomial defined in the support \CWthm, $\prod_{i=0}^{n-1}\sum_{j=1}^n\alpha_j^{q^i}X_j$ where $\{\alpha_j\colon j\in[n]\}$ was a basis of $\bbF_{q^n}$ over $\bbF_q$. But in this example we take the same polynomial but on $d$ variables but the solution space is over $\bbF_{q^n}$. Hence our polynomial in this case is $f(X_1,\dots,X_n)=g(X_1,\dots,X_d)=\prod_{i=0}^{n-1}\sum_{j=1}^d\alpha_j^{q^i}X_j$. Since $f$ had only one solution, $g$ has a solution at $(\gm_1,\dots,\gm_n)$ if $\gm_i=0$ for all $i\in[d]$. Therefore $g$ has exactly $q^{n-d}$ solutions in $\bbF_q^n$. 

\subsection{Upper Bounds on the Number of Solutions}

We now give upper bounds on the number of solutions of a polynomial equation. The first is the celebrated Schwartz--Zippel lemma, which bounds the probability of a random point being a root in terms of total degree.
\begin{lemma}{Schwartz-Zippel Lemma}{lem:schwartz-zippel}
  Let $f\in\bbF_q[X_1,\dots, X_n]$ with $\deg(f)=d$ be a non-zero polynomial. If $S\subseteq \bbF_q$ be any subset then $$\Pr[\alpha_1,\dots,\alpha_n\in S] [f(\alpha_1,\dots,\alpha_n)=0]\leq \frac{d}{|S|}$$
\end{lemma}
This lemma immediately gives an upper bound of $d\cdot q^{n-1}$ solutions for the equations $f(X_1,\dots, X_n)=0$ in $\bbF_q^n$. Hence 
\begin{corolary}{}{}
  Let $f\in\bbF_q[X_1,\dots, X_n]$ with $\deg(f)=d$ be a non-zero polynomial. Then the equation $f(X_1,\dots, X_n)=0$ has at most $d\cdot q^{n-1}$ solutions in $\bbF_q^n$. 
\end{corolary}

We can give another probabilistic upper bound when we control the degree of $f$ in each variable separately rather than just the total degree.
\begin{Theorem}{Individual Degree Bound}{thm:individual-degree-bound}
  Let $f\in\bbF_q[X_1,\dots,X_n]$ be a non-zero polynomial with $\deg_{X_i}(f)\leq d$ for all $i\in[n]$. Then for $(\alpha_1,\dots,\alpha_n)$ chosen uniformly at random from $\bbF_q^n$,
  $$\Pr[\alpha_1,\dots,\alpha_n\in\bbF_q][f(\alpha_1,\dots,\alpha_n)=0]\leq 1-\lt(1-\frac{d}{q}\rt)^{n}$$
\end{Theorem}
\begin{proof}
  We proceed by induction on $n$. For the base case $n=1$, since $f\in\bbF_q[X_1]$ is non-zero with $\deg_{X_1}(f)\leq d$, it has at most $d$ roots in $\bbF_q$. Hence $\Pr[x_1\in\bbF_q][f(x_1)=0]\leq d/q=1-(1-d/q)^1$. 
  
  Assume the bound holds for $n-1$ variables. Let $t=\deg_{X_n}(f)$. Then write $f$ as a polynomial in $X_n$ with coefficients in $\bbF_q[X_1,\dots,X_{n-1}]$:
  $$f(X_1,\dots,X_n)=\sum_{i=0}^{t}f_i(X_1,\dots,X_{n-1})\cdot X_n^k.$$
  Since $\deg_{X_j}(f)\leq d$ for each $j\in[n-1]$, every $f_i$ satisfies $\deg_{X_j}(f_i)\leq d$.

  Let $\alpha_1,\dots,\alpha_{n-1}$ be uniformly random points on $\bbF_q$ and $\alpha_n$ be uniform on $\bbF_q$, independently. Then by induction hypothesis $$\Pr[\alpha_1,\dots,\alpha_{n-1}\in\bbF_q] [f_t(\alpha_1,\dots,\alpha_{n-1})=0]\leq 1-\lt(1-\frac{d}{q}\rt)^{n-1}$$Now we have \begin{align*}
    \underset{\alpha_1,\dots,\alpha_n\in\bbF_q}{\bbP} [f(\alpha_1,\dots,\alpha_n)=0] &\leq \underset{\alpha_1,\dots,\alpha_{n-1}\in\bbF_q}{\bbP} [f_t(\alpha_1,\dots,\alpha_{n-1})=0]+ \underset{\alpha_1,\dots,\alpha_n\in\bbF_q}{\bbP} [f(\alpha_1,\dots, \alpha_{n-1},\alpha_n)=0\wedge f_t(\alpha_1,\dots,\alpha_{n-1})\neq 0]\\ 
    & \leq \underset{\alpha_1,\dots,\alpha_{n-1}\in\bbF_q}{\bbP} [f_t(\alpha_1,\dots,\alpha_{n-1})=0]\;+ \\ 
    &\qquad\qquad \underset{\alpha_n\in\bbF_q}{\bbP} [f(\alpha_1,\dots, \alpha_{n-1},\alpha_n)=0\mid f_t(\alpha_1,\dots,\alpha_{n-1})\neq 0]\underset{\alpha_1,\dots,\alpha_{n-1}\in\bbF_q}{\bbP}[f_t(\alpha_1,\dots,\alpha_{n-1})\neq 0]
  \end{align*}
  Let $p\coloneqq \underset{\alpha_1,\dots,\alpha_{n-1}\in\bbF_q}{\bbP} [f_t(\alpha_1,\dots,\alpha_{n-1})=0]$. For the second term after fixing $\alpha_1,\dots,\alpha_{n-1}$ we have $\underset{\alpha_n\in\bbF_q}{\bbP} [f(\alpha_1,\dots, \alpha_{n-1},\alpha_n)=0\mid f_t(\alpha_1,\dots,\alpha_{n-1})\neq 0]\leq\frac{d}q$. Now $$\underset{\alpha_1,\dots,\alpha_{n-1}\in\bbF_q}{\bbP}[f_t(\alpha_1,\dots,\alpha_{n-1})\neq 0]=1-\underset{\alpha_1,\dots,\alpha_{n-1}\in\bbF_q}{\bbP}[f_t(\alpha_1,\dots,\alpha_{n-1})= 0]$$ Therefore together we get $$\underset{\alpha_1,\dots,\alpha_n\in\bbF_q}{\bbP} [f(\alpha_1,\dots,\alpha_n)=0]\leq p+(1-p)\frac{d}{q}=\frac{d}{q}+\lt(1-\frac{d}{q}\rt)p$$By induction hypothesis $p\leq 1-\lt(1-\frac{d}q\rt)^{n-1}$. Therefore we have $$\underset{\alpha_1,\dots,\alpha_n\in\bbF_q}{\bbP} [f(\alpha_1,\dots,\alpha_n)=0]\leq \frac{d}{q}+\lt(1-\frac{d}{q}\rt)\lt[1-\lt(1-\frac{d}{q}\rt)^{n-1}\rt]=1-\lt(1-\frac{d}{q}\rt)^{n}$$So we have the lemma. 
\end{proof}

We can give another upper bound based on individual degrees of each variables. If degree of each variable is at most $d$ then there can be at most $d(q^{n-1}-1)$ many non-trivial solutions. For this we consider $f$ to be \emph{homogeneous} i.e. total degree of each monomial is same. 
\begin{Theorem}{}{}
  Let $f\in\bbF_q[X_1,\dots, X_n]$ be a homogeneous with $\deg(f)=d\geq 1$. Then the equation $f(X_1,\dots,X_n)=0$ has at most $d(q^{n-1}-1)$ non-trivial solutions. 
\end{Theorem}
\begin{proof}
  We use double induction on $n$ and $d$. For the base case $n=1$ then $f=aX_1^d$ for some $a\neq 0$, so the only root is $X_1=0$, giving $\ov{Z}(f)=0\leq d(q^0-1)=0$. If $d=1$ then $f=a_1X_1+\cdots+a_nX_n$ is a non-zero linear form whose zero set is an $(n-1)$-dimensional hyperplane, so $\ov{Z}(f)=q^{n-1}-1=d(q^{n-1}-1)$.

  Suppose $n>1$, $d>1$, and the bound holds for non-constant homogeneous polynomials in at most $n$ variables of degree less than $d$ and for those in fewer than $n$ variables of degree at most $d$. We consider two cases.

  \paragraph*{Case 1: $X_1\mid f$.} Write $f=X_1\cdot g$ where $g\in\bbF_q[X_1,\dots,X_n]$ is homogeneous of degree $d-1$. Every nontrivial solution $(\alpha_1,\dots,\alpha_n)\in\ovZ(f)$ with $\alpha_i\in\bbF_q$ for all $i\in[n]$, satisfies $\alpha_1=0$ or $g(\alpha_1,\dots,\alpha_n)=0$. Solutions with $\alpha_1=0$: any $(0,\alpha_2,\dots,\alpha_n)\neq 0$ works, giving $q^{n-1}-1$ solutions. Solutions with $\alpha_1\neq 0$: these are nontrivial zeros of $g$, and by the induction hypothesis on $d$ there are at most $(d-1)(q^{n-1}-1)$. Hence $$\ov{Z}(f)\leq (q^{n-1}-1)+(d-1)(q^{n-1}-1)=d(q^{n-1}-1).$$

  \paragraph*{Case 2: $X_1\nmid f$.} Write $f=\sum_{k=0}^{d}f_k(X_2,\dots,X_n)\cdot X_1^k$ with $f_k$ homogeneous of degree $d-k$. Since $X_1\nmid f$ we have $f_0\neq 0$, so for every $\alpha\in\bbF_q$ such that $f(\alpha,X_2,\dots,X_n)$ is a non-zero polynomial of degree $d$ in $n-1$ variables we have by \SZlem it has at most $d\cdot q^{n-2}$ zeros in $\bbF_q^{n-1}$. Now summing over $\alpha\in\bbF_q^*$ gives at most $(q-1)\cdot d\cdot q^{n-2}$ nontrivial solutions with $\alpha_1\neq 0$. For solutions with $\alpha_1=0$: $f(0,X_2,\dots,X_n)=f_0(X_2,\dots,X_n)$ is a nonzero homogeneous polynomial in $n-1$ variables of degree $d$, so by the induction hypothesis on $n$ it has at most $d(q^{n-2}-1)$ nontrivial zeros. Altogether, $\ov{Z}(f)\leq (q-1)d\cdot q^{n-2}+d(q^{n-2}-1)=d\lt(q^{n-1}-q^{n-2}+q^{n-2}-1\rt)=d(q^{n-1}-1)$.
\end{proof}
\subsection{Average Number of Solutions}
We now investigate the \emph{average number of solutions} of a polynomial equations. Let $d\in\bbN$. Then define $\Om_d\coloneqq \bbF_q^{\leq d}[X_1,\dots, X_n]$ and $\om(d)$ be the set of $n-$tuples $(i_1,\dots,i_n)\in\bbZ_0$ such that $i_1+\cdots +i_n\leq d$. Then 
\begin{observation}\label{obs:Omd-omd}
  With the definitions of $\Om_d$ and $\om(d)$ as above we have $|\Om_d|=q^{|\om(d)|}$. 
\end{observation}
\begin{lemma}{}{lem:average-solutions}
  With the notations of $\Om_d$ and $\om(d)$ as defined above we have $$\frac1{|\Om_d|}\sum_{f\in\Om_d}Z(f)=q^{n-1}$$
\end{lemma}
\begin{proof}
  We have $$\sum_{f\in\Om_d}Z(f)=\sum_{f\in\Om_d}\sum_{\substack{\alpha_1,\dots,\alpha_n\in\bbF_q\\ f(\alpha_1,\dots,\alpha_n)=0}}1=\sum_{\alpha_1,\dots,\alpha_n\in\bbF_q}\sum_{\substack{f\in\Om_d\\ f(\alpha_1,\dots,\alpha_n)=0}}1$$Let for any $f\in\Om_d$ $$f(X_1,\dots,X_n)=\sum_{(i_1,\dots,i_n)\in\om(d)}f_{i_1,\dots,i_n}X_1^{i_1}\cdots X_n^{i_n}\quad f_{i_1,\dots,i_n}\in\bbF_q,\ \forall\ (i_1,\dots, i_n)\in\om(d)$$Then for a fixed $\alpha_1,\dots, \alpha_n\in\bbF_q$ we get to choose the coefficients $f_{i_1,\dots, i_n}$ for all $(i_1,\dots,i_n)\in\om(d)\setminus\{(0,\dots, 0)\}$ arbitrarily and then determining the constant term $f_{0,\dots, o}$ to make $f(\alpha_1,\dots,\alpha_n)=0$. Thus the number of $f\in\om_d$ with $f(\alpha_1,\dots, \alpha_n)=0$ is equal to $q^{|\om(d)|-1}$. Therefore $\sum_{f\in\Om_d}Z(f)=q^n\cdot q^{|\om(d)|-1}=|\Om_d|\cdot q^{n-1}$. So we have the lemma. 
\end{proof}
The proof of the above lemma gives us the following observation:
\begin{observation}\label{obs:num-poly-zero-at-point}
  For any $\alpha_1,\dots,\alpha_n\in\bbF_q$, the number of polynomials $f\in\Om_d$ such that $f(\alpha_1,\dots,\alpha_n)=0$ is $q^{|\om(d)|-1}$. 
\end{observation}

The above lemma implies that a polynomial equation in $n$ variables has on average $q^{n-1}$ solutions in $\bbF_q^n$. So we now consider average deviation from the average number of solutions, $q^{n-1}$. 
\begin{Theorem}{}{}
  With the notations of $\Om_d$ and $\om(d)$ as defined above we have
  $$\frac1{|\Om_d|}\sum_{f\in\Om_d}\Big(Z(f)-q^{n-1}\Big)^2=q^{n-1}-q^{n-2}$$
\end{Theorem}
\begin{proof}
  Let $\alpha=(\alpha_1,\dots,\alpha_n),\beta=(\beta_1,\dots,\beta_n)\in\bbF_q^n$ be any two points. Then we get
  $$\sum_{f\in\Om_d}Z(f)^2  = \sum_{f\in\Om_d}\lt( \sum_{\substack{\alpha\in\bbF_q^n\\ f(\alpha)=0}} 1 \rt)^2= \sum_{f\in\Om_d}\sum_{\substack{\alpha\in\bbF_q^n\\ f(\alpha)=0}}\sum_{\substack{\beta\in\bbF_q^n\\ f(\beta)=0}}1=\sum_{\alpha,\beta\in\bbF_q^n}\sum_{\substack{f\in\Om_d\\  f(\alpha)=f(\beta)=0}}1$$
  Now if $\alpha=\beta$ then by \obsref{obs:num-poly-zero-at-point} we have $\sum_{f\in\Om_d,\   f(\alpha)=0}1=q^{|\om(d)|-1}$. Now suppose $\alpha\neq\beta$. This gives two linear equations for the coefficients of $f$. Therefore there are $q^{|\om(d)|-2}$ polynomials $f\in\Om_d$ such that $f(\alpha)=f(\beta)=0$. Hence we obtain \begin{align*}
    \sum_{f\in\Om_d}Z(f)^2  =  & = \sum_{\alpha\in\bbF_q^n}|\Om_d|\cdot q^{-1}+\sum_{\alpha,\beta\in\bbF_q^n,\alpha\neq \beta}|\Om_d|\cdot q^{-2} \byo{obs:Omd-omd}\\ 
    & = q^n\cdot |\Om_d|\cdot q^{-1}+q^n(q^n-1)|\Om_d|\cdot q^{-2}\\ 
    & = |\Om_d|\Big( q^{n-1}+q^{2n-2}-q^{n-2} \Big)
  \end{align*}
  So now we get \begin{align*}
    \sum_{f\in\Om_d}(Z(f)-q^{n-1})^2 & = \sum_{f\in\Om_d}Z(f)^2 + 2\cdot q^{n-1}\sum_{f\in\Om_d}Z(f)+q^{2n-2}\sum_{f\in\Om_d}1\\ 
    & = |\Om_d|\Big(q^{2n-2}+q^{n-1}-q^{n-2}\Big)-2\cdot q^{n-1}\cdot |\Om_d|\cdot q^{n-1}+q^{2n-2}\cdot |\Om_d|\\ 
    & = |\Om_d|(q^{n-1}-q^{n-2})
  \end{align*}
  So we have the result. 
\end{proof}
From this result we can expect that $|Z(f)-q^{n-1}|$ is often $O\lt(q^{\nicefrac{(q-1)}{2}}\rt)$. Later we will see instances of such expected behavior for various polynomial equations. In the next sections we will talk about character sums over polynomial evaluations. 

\section{Character Sums with Polynomials}
\section{Average Number of Solutions}\label{sec:average-solutions}
We now investigate the \emph{average number of solutions}\index{average number of solutions} of a polynomial equations. Let $d\in\bbN$. Then define $\Om_d\coloneqq \bbF_q^{\leq d}[X_1,\dots, X_n]$ and $\om(d)$ be the set of $n-$tuples $(i_1,\dots,i_n)\in\bbZ_0$ such that $i_1+\cdots +i_n\leq d$. Then 
\begin{observation}\label{obs:Omd-omd}
  With the definitions of $\Om_d$ and $\om(d)$ as above we have $|\Om_d|=q^{|\om(d)|}$. 
\end{observation}
\begin{lemma}{Average of Solution Count}{lem:average-solutions}
  \index{average number of solutions}%
  With the notations of $\Om_d$ and $\om(d)$ as defined above we have $$\frac1{|\Om_d|}\sum_{f\in\Om_d}Z(f)=q^{n-1}$$
\end{lemma}
\begin{proof}
  We have $$\sum_{f\in\Om_d}Z(f)=\sum_{f\in\Om_d}\sum_{\substack{\alpha_1,\dots,\alpha_n\in\bbF_q\\ f(\alpha_1,\dots,\alpha_n)=0}}1=\sum_{\alpha_1,\dots,\alpha_n\in\bbF_q}\sum_{\substack{f\in\Om_d\\ f(\alpha_1,\dots,\alpha_n)=0}}1$$Let for any $f\in\Om_d$ $$f(X_1,\dots,X_n)=\sum_{(i_1,\dots,i_n)\in\om(d)}f_{i_1,\dots,i_n}X_1^{i_1}\cdots X_n^{i_n}\quad f_{i_1,\dots,i_n}\in\bbF_q,\ \forall\ (i_1,\dots, i_n)\in\om(d)$$Then for a fixed $\alpha_1,\dots, \alpha_n\in\bbF_q$ we get to choose the coefficients $f_{i_1,\dots, i_n}$ for all $(i_1,\dots,i_n)\in\om(d)\setminus\{(0,\dots, 0)\}$ arbitrarily and then determining the constant term $f_{0,\dots, o}$ to make $f(\alpha_1,\dots,\alpha_n)=0$. Thus the number of $f\in\om_d$ with $f(\alpha_1,\dots, \alpha_n)=0$ is equal to $q^{|\om(d)|-1}$. Therefore $\sum_{f\in\Om_d}Z(f)=q^n\cdot q^{|\om(d)|-1}=|\Om_d|\cdot q^{n-1}$. So we have the lemma. 
\end{proof}
The proof of the above lemma gives us the following observation:
\begin{observation}\label{obs:num-poly-zero-at-point}
  For any $\alpha_1,\dots,\alpha_n\in\bbF_q$, the number of polynomials $f\in\Om_d$ such that $f(\alpha_1,\dots,\alpha_n)=0$ is $q^{|\om(d)|-1}$. 
\end{observation}

The above lemma implies that a polynomial equation in $n$ variables has on average $q^{n-1}$ solutions in $\bbF_q^n$. So we now consider average deviation from the average number of solutions, $q^{n-1}$. 
\begin{Theorem}{Variance of Solution Count}{}
  \index{solution count!variance}%
  With the notations of $\Om_d$ and $\om(d)$ as defined above we have
  $$\frac1{|\Om_d|}\sum_{f\in\Om_d}\Big(Z(f)-q^{n-1}\Big)^2=q^{n-1}-q^{n-2}$$
\end{Theorem}
\begin{proof}
  Let $\alpha=(\alpha_1,\dots,\alpha_n),\beta=(\beta_1,\dots,\beta_n)\in\bbF_q^n$ be any two points. Then we get
  $$\sum_{f\in\Om_d}Z(f)^2  = \sum_{f\in\Om_d}\lt( \sum_{\substack{\alpha\in\bbF_q^n\\ f(\alpha)=0}} 1 \rt)^2= \sum_{f\in\Om_d}\sum_{\substack{\alpha\in\bbF_q^n\\ f(\alpha)=0}}\sum_{\substack{\beta\in\bbF_q^n\\ f(\beta)=0}}1=\sum_{\alpha,\beta\in\bbF_q^n}\sum_{\substack{f\in\Om_d\\  f(\alpha)=f(\beta)=0}}1$$
  Now if $\alpha=\beta$ then by \obsref{obs:num-poly-zero-at-point} we have $\sum_{f\in\Om_d,\   f(\alpha)=0}1=q^{|\om(d)|-1}$. Now suppose $\alpha\neq\beta$. This gives two linear equations for the coefficients of $f$. Therefore there are $q^{|\om(d)|-2}$ polynomials $f\in\Om_d$ such that $f(\alpha)=f(\beta)=0$. Hence we obtain \begin{align*}
    \sum_{f\in\Om_d}Z(f)^2  =  & = \sum_{\alpha\in\bbF_q^n}|\Om_d|\cdot q^{-1}+\sum_{\alpha,\beta\in\bbF_q^n,\alpha\neq \beta}|\Om_d|\cdot q^{-2} \byo{Omd-omd}\\ 
    & = q^n\cdot |\Om_d|\cdot q^{-1}+q^n(q^n-1)|\Om_d|\cdot q^{-2}\\ 
    & = |\Om_d|\Big( q^{n-1}+q^{2n-2}-q^{n-2} \Big)
  \end{align*}
  So now we get \begin{align*}
    \sum_{f\in\Om_d}(Z(f)-q^{n-1})^2 & = \sum_{f\in\Om_d}Z(f)^2 + 2\cdot q^{n-1}\sum_{f\in\Om_d}Z(f)+q^{2n-2}\sum_{f\in\Om_d}1\\ 
    & = |\Om_d|\Big(q^{2n-2}+q^{n-1}-q^{n-2}\Big)-2\cdot q^{n-1}\cdot |\Om_d|\cdot q^{n-1}+q^{2n-2}\cdot |\Om_d|\\ 
    & = |\Om_d|(q^{n-1}-q^{n-2})
  \end{align*}
  So we have the result. 
\end{proof}
From this result we can expect that $|Z(f)-q^{n-1}|$ is often $O\lt(q^{\nicefrac{(q-1)}{2}}\rt)$. Later we will see instances of such expected behavior for various polynomial equations. In the next sections we will talk about character sums over polynomial evaluations. 

\section{Quadratic Forms}\label{sec:quadratic-forms}
In this section we will study about quadratic forms and its equivalent forms which are easier to study than a general quadratic form. We will give bounds on the deviation of number of solutions from average number of solutions. Before all of that let's see what is a quadratic form.
\begin{Definition}{Quadratic Forms}{}
  \index{quadratic form}%
  A quadratic form over $\bbF_q$ is a degree-$2$ \emph{homogeneous} polynomial or the zero polynomial.
\end{Definition}
If $f\in\bbF_q[X_1,\dots,X_n]$ is a quadratic form. Now if $q$ is odd then for any $i,j\in[n]$ we write $a_{i,j}X_iX_j=\frac12a_{i,j}X_iX_j+\frac12a_{ij}X_jX_i$ which leads to the representation $$f(X_1,\dots, X_n)= \sum_{i,j\in[n]}a_{i,j}X_iX_j\quad\text{for all }a_{i,j}\in\bbF_q\text{ with }a_{i,j}=a_{j,i}.$$
From this formulation we can represent $f$ by a matrix denoted by $C_f$ where $(i,j)^{th}$ entry of $C_f$ is $a_{i,j}$. This matrix is called the \emph{coefficient matrix}\index{coefficient matrix}.
\begin{observation}\label{obs:coeff-matrix-symmetric}
  Since $a_{i,j}=a_{j,i}$ we have $C_f^\top=C_f$.
\end{observation}
So if $X$ denotes the column vector of all variables then $f=X^\top C_fX$. With this representation now we can talk about equivalence of two quadratic forms.
\begin{Definition}{Equivalent Quadratic Forms}{}
  \index{quadratic form!equivalent}%
  Let $f,g\in\bbF_q[X_1,\dots, X_n]$ be two quadratic forms over $\bbF_q$. Then $f$ and $g$ are said to be \emph{equivalent} if $f$ can be transformed into $g$ by a non-singular linear transformation of the variables.
\end{Definition}

So let $f,g\in\bbF_q[X_1,\dots, X_n]$ are two equivalent quadratic forms. Then Let $X$ be the variable vector for $f$ and $\tdX$ be the variable vector for $g$. Then $f=X^\top C_fX$ and $g=\tdX^\top C_g\tdX$. Suppose $M$ is the non-singular linear transformation of variables i.e.\ $\tdX=MX$. For odd $q$ we have $$\tdX^\top C_g\tdX=g=(MX)^\top C_f(MX)= X^\top (M^\top C_f M)X,$$so $C_g=M^\top C_fM$.
\begin{observation}\label{obs:coeff-matrix-equivalent}
  If $f,g\in\bbF_q[X_1,\dots, X_n]$ ($q$ odd) are two quadratic forms then $f$ and $g$ are equivalent if and only if there exists an invertible matrix $M\in\bbF_q^{n\times n}$ such that $C_g=M^\top C_fM$ where $M$ is the non-singular linear transformation. 
\end{observation}
Since $\tdX=MX$. So $M$ gives the one-one correspondence between $\mcV(f(X_1,\dots,X_n)=\alpha)$ and $\mcV(g(X_1,\dots,X_n)=\alpha)$ for any $\alpha\in\bbF_q$. 
\begin{observation}\label{obs:solution-correspondence-equivalent}
  If $f,g\in\bbF_q[X_1,\dots, X_n]$ ($q$ odd) are two quadratic forms then $f$ and $g$ are equivalent by the non-singular linear transformation $M\in\bbF_q^{n\times n}$. Then $\gm\mapsto M\gm$ where $\gm\in\bbF_q^n$ gives an one-one correspondence between $\mcV(f(X_1,\dots,X_n)=\alpha)$ and $\mcV(g(X_1,\dots,X_n)=\alpha)$ for any $\alpha\in\bbF_q$. 
\end{observation}
We will use another terminology to make our life easier. If $f\in\bbF_q[X_1,\dots, X_n]$ is a quadratic form then for any $\alpha\in\bbF_q$ we will say $f$ \emph{represents} $\alpha$ if $f(X_1,\dots, X_n)=\alpha$ has a solution.

To calculate and understand the of solutions of quadratic forms we define the integer-valued function $\vth:\bbF_q\to\bbZ$ by $\vth(\alpha)=-1$ for all $\alpha\in\bbF_q^*$ and $\vth(0)=q-1$. Then we have the following properties of $\vth$ function.\index{theta-function@$\theta$-function}
\begin{lemma}{}{lem:theta-convolution}
  For any finite field $\bbF_q$ we have $$\sum_{\alpha\in\bbF_q}\vth(\alpha)=0,$$and for any $\alpha\in\bbF_q$ and integers $1\leq k\leq m$,
  $$\sum_{\substack{\alpha_1,\dots,\alpha_m\in\bbF_q\\ \alpha_1+\cdots+\alpha_m=\alpha}}\vth(\alpha_1)\cdots \vth(\alpha_k)=\begin{cases}
    0 & \text{if $1\leq k<m$}\\
    \vth(\alpha)\cdot q^{m-1} & \text{if $k=m$.}
  \end{cases}$$
\end{lemma}
\begin{proof}
  Since for all $\alpha\in\bbF_q^*$, $\vth(\alpha)=-1$ we have $\sum_{\alpha\in\bbF_q^*}\vth(\alpha)=-(q-1)$. Therefore $\vth(0)+\sum_{\alpha\in\bbF_q^*}\vth(\alpha)=0$. So we have the first result. 

  Now suppose $k<m$. Then we have \begin{align*}
    \sum_{\substack{\alpha_1,\dots,\alpha_m\in\bbF_q\\ \alpha_1+\cdots+\alpha_m=\alpha}}\prod_{j=1}^k\vth(\alpha_j) & = \sum_{\alpha_1,\dots,\alpha_k\in\bbF_q}\lt(\prod_{j=1}^k\vth(\alpha_j)\rt)\sum_{\substack{\alpha_{k+1},\dots,\alpha_m\in\bbF_q\\ \alpha_{k+1}+\cdots+\alpha_m=\alpha-\alpha_1-\cdots-\alpha_k}}1\\ 
    & = q^{m-k-1}\sum_{\alpha_1,\dots,\alpha_k\in\bbF_q}\prod_{j=1}^k\vth(\alpha_j)\\ 
    & = q^{m-k-q}\prod_{j=1}^k\sum_{\alpha_j\in\bbF_q}\vth(\alpha_j) =0
  \end{align*}
  Now assume $k=m$. Here we will use induction. The case $m=1$ holds trivially. Suppose the formula is shown for $m\geq 1$. Then \begin{align*}
    \sum_{\substack{\alpha_1,\dots,\alpha_{m+1}\in\bbF_q\\ \alpha_1+\cdots+\alpha_{m+1}=\alpha}}\prod_{j=1}^{m+1}\vth(\alpha_j) & = \sum_{\substack{\alpha_1,\dots,\alpha_{m+1}\in\bbF_q\\ \alpha_1+\cdots+\alpha_{m+1}=\alpha}}\prod_{j=1}^{m+1}\vth(\alpha_j)+\sum_{\substack{\alpha_1,\dots,\alpha_{m+1}\in\bbF_q\\ \alpha_1+\cdots+\alpha_{m+1}=\alpha}}\prod_{j=1}^m\vth(\alpha_j)\\ 
    & = \sum_{\substack{\alpha_1,\dots,\alpha_{m+1}\in\bbF_q\\ \alpha_1+\cdots+\alpha_{m+1}=\alpha}}\lt[\prod_{j=1}^m\vth(\alpha_j)\rt]\cdot [\vth(\alpha_{m+1})+1]\\ 
    & = \sum_{\alpha_1,\dots,\alpha_{m}\in\bbF_q}\lt[\prod_{j=1}^m\vth(\alpha_j)\rt]\cdot \lt[\vth\lt(\alpha-\sum_{j=1}^m\alpha_j\rt)+1\rt]\\ 
    & = q\sum_{\substack{\alpha_1,\dots,\alpha_{m}\in\bbF_q\\ \alpha_1+\cdots+\alpha_{m}=\alpha}}\prod_{j=1}^m\vth(\alpha_j)& [\text{If $\alpha_1+\cdots+\alpha_m\neq \alpha$ then the term is $0$}]\\ 
    & = q\cdot q^{m-2}\cdot \vth(\alpha) & [\text{By Induction Hypothesis}]\\ 
    & = q^{m-1}\cdot \vth(\alpha)
  \end{align*}
  so now we have the second result too. So we have the lemma. 
\end{proof}
\subsection{Odd Characteristic Quadratic Forms}
\index{quadratic form!odd characteristic}%
We now resume the study of quadratic forms $f\in\bbF_q[X_1,\dots,X_n]$ where $q$ is odd. We will first show that every quadratic form $f\in\bbF_q[X_1,\dots,X_n]$ where $q$ is odd is equivalent to a \emph{diagonal quadratic form}\index{quadratic form!diagonal} i.e. a quadratic form of type $a_1X_q^2+\cdots+a_nX_n^2$ where $a_i\in\bbF_q$ for all $i\in[n]$. 
\begin{lemma}{Reduction by Separating a Variable}{lem:qf-reduction-by-rep}
  Let $q$ be odd and $f\in\bbF_q[X_1,\dots,X_n]$ be a quadratic form with $n\geq 2$. If $f$ represents $\alpha\in\bbF_q^*$ then $f$ is equivalent to $\alpha\cdot X_1^2+g(X_2,\dots, X_n)$ where $g$ is a quadratic form in $n-1$ variables.
\end{lemma}
\begin{proof}
  Since $f$ represents $\alpha$, $\exs\ (\gm_1,\dots,\gm_n)\in\bbF_q^n$ such that $f(\gm_1,\dots, \gm_n)=\alpha$. Since $\alpha\neq 0$ not all $\gm_i$ are zero. So we can find an invertible matrix $M\in\bbF_q^{n\times n}$ such that in the first column of $M$ the entries are $\gm_1,\dots,\gm_n$. Therefore if we apply the non-singular linear transformation by $M$ to $f$, we obtain a quadratic form in $Y_1,\dots, Y_n$ for which the coefficient of $Y_1^2$ is $f(\gm_1,\dots, \gm_n)=\alpha$. Thus $f$ is equivalent to a quadratic form of the type $$\alpha Y_1^2+2b_2Y_1Y_2+\cdots +2bY_1Y_n+h(Y_2,\dots, Y_n)=\alpha(Y_1+b_2\cdot \alpha^{-1}\cdot Y_2+\cdots b_n\cdot \alpha^{-1}\cdot Y_n)^2+g(Y_2,\dots, Y_n)$$ for some $b_2,\dots, b_n\in\bbF_q$ where $g,h\in\bbF_q[X_1,\dots,X_n]$ are two quadratic forms. So now take the non-singular linear transformation \begin{align*}
    Z_1 & = Y_1+b_2\cdot \alpha^{-1}\cdot Y_2+\cdots +b_n\cdot \alpha^{-1}\cdot Y_n&&\\ 
    Z_i & = Y_i\ \forall\ i\in\{2,\dots, n\}&&
  \end{align*}Then we get a quadratic form of the desired type and this is equivalent to $f$. 
\end{proof}

Applying \QFRedlem repeatedly separates variables one by one, eventually reaching a diagonal form. The only subtlety is ensuring at each step that the remaining form represents a non-zero element of $\bbF_q$; this is handled in the proof of \QFDiagthm below.


\begin{Theorem}{Diagonal Form of Quadratic Forms over $\bbF_q$}{thm:qf-diagonal-odd}
  For any quadratic form $f\in\bbF_q[X_1,\dots,X_n]$ where $q$ is odd, there exists $a_1,\dots, a_n\in\bbF_q$ such that $f$ is equivalent to the diagonal quadratic form $a_1 X_1^2+\cdots +a_nX_n^2$.
\end{Theorem}
\begin{proof}
  We will prove this using induction on the number of variables. If $n=1$ then $f(X_1)=a_{1,1}X_1^2$. Then its already diagonal. So suppose $n\geq 2$ and the result holds for quadratic forms in $(n-1)$ variables. 
  
  Let $f\in\bbF_q[X_1,\dots,X_n]$ be a quadratic form in $n$ variables. Now the theorem is trivial if $f$ is a zero polynomial. So suppose $f$ is a non-zero polynomial. If there exists $i\in[n]$, $a_{i,i}\neq 0$ then $f$ represents $a_{i,i}$ since at the point $X_i=1$ and $X_j=0$ for all $j\in[n]\setminus\{i\}$, $f$ evaluates to be $a_{i,i}$. If for all $i\in[n]$, $a_{i,i}=0$ then $\exs\ i,j\in[n]$, $i\neq j$ such that $a_{i,j}=a_{j,i}\neq 0$. Then $f$ represents $2a_{i,j}$ since at the point $X_i=X_j=1$ and $X_t=0$ for all $t\in[n]\setminus\{i,j\}$, $f$ evaluates to be $2a_{i,j}$. 
  
  So $f$ represents some element $a_1\in\bbF_q^*$ and hence by \QFRedlem we have $f$ is equivalent to a quadratic form $a_1 X_1^2+g(X_2,\dots, X_n)$ where $g$ is a quadratic form in $(n-1)$ variables. Now by induction hypothesis $g$ is equivalent to a diagonal quadratic form $a_2X_2^2+\cdots a_nX_n^2$. Therefore $f$ is equivalent to the diagonal quadratic form $a_1 X_1^2+\cdots +a_nX_n^2$. 
\end{proof}

So now we can only talk about the solutions of diagonal quadratic forms since all general quadratic forms in $\bbF_q[X_1,\dots,X_n]$ are equivalent to a diagonal quadratic form. Now let $f\in\bbF_q[X_1,\dots,X_n]$ is equivalent to $a_1X_1^2+\cdots+a_nX_n^2$. Then some $a_i$'s may be zero. Since non-singular linear transformation preserves ranks, equivalent quadratic forms have coefficient matrices with same rank. Therefore non-singular linear transformation also preserves the non-zero ness of determinant of the coefficient matrix. 
\begin{observation}\label{obs:diagonal-rank-det}
  Let $f\in\bbF_q[X_1,\dots,X_n]$ is a quadratic form is equivalent to $a_1X_1^2+\cdots+a_nX_n^2$. Then $\Rank(C_f)=n-|\{i\in[n]\colon a_i=0\}|$. Similarly, $\det(C_f)=0$ if and only if $\prod_{j=1}^na_i=0$. 
\end{observation}
\begin{observation}\label{obs:equivalent-det-value}
  Let $f,g\in\bbF_q[X_1,\dots,X_n]$ are two equivalent quadratic forms by the non-singular linear transformation $M\in\bbF_q^{n\times n}$. Since by \QFEquivCoeffObs we have $C_g=M^\top C_f M$ and therefore $\det(g)=\det(f)\cdot \det(M)^2$. 
\end{observation}
If the coefficient matrix $C_f$ for any quadratic form $f\in\bbF_q[X_1,\dots,X_n]$ has rank $n$ then we call $f$ to be \emph{non-degenerate}\index{quadratic form!non-degenerate}. We also define \emph{determinant of $f$}\index{quadratic form!determinant}, $\det(f)$ to be the determinant of the coefficient matrix $C_f$. 

The two-variable lemma below is the base case for the induction used in \QFCountOddthm. Once we know $Z(a_1X_1^2+a_2X_2^2=\alpha)$ exactly, the convolution property \QFConvProplem allows us to build up the solution count for $n$ variables.

\begin{lemma}{Solutions of a Two-Variable Diagonal Quadratic Form}{lem:qf-two-var-count}
  For odd $q$, let $\alpha\in\bbF_q$, $a_1,a_2\in\bbF_q^*$, and let $\eta\in\Mq$ be the quadratic character of $\bbF_q$. Then $$Z\lt(a_1X_1^2+a_2X_2^2=\alpha\rt)=q+\vth(\alpha)\cdot \eta(-a_1a_2).$$
\end{lemma}
\begin{proof}
  Now we can break $Z(a_1X_1^2+a_2X_2^2=\alpha)$ into sum over all $\alpha_1,\alpha_2\in\bbF_q$ product of $Z(a_1X_1^2=\alpha_1)$ and $Z(a_2X_2^2=\alpha_2$. So we have 
  \begin{align*}
    Z(a_1X_1^2+a_2X_2^2=\alpha) & = \sum_{\substack{\alpha_1,\alpha_2\in\bbF_q\\ \alpha_1+\alpha_2=\alpha}}Z(a_1X_1^2=\alpha_1)\cdot Z(a_2X_2^2=\alpha_2)\\ 
    & = \sum_{\substack{\alpha_1,\alpha_2\in\bbF_q\\ \alpha_1+\alpha_2=\alpha}} (1+\eta(a_1^{-1}\alpha_1))\cdot (1+\eta(a_2^{-1}\alpha_2))\\ 
    & = \sum_{\substack{\alpha_1,\alpha_2\in\bbF_q\\ \alpha_1+\alpha_2=\alpha}} \lt(1+\eta(a_1^{-1}\alpha_1)+\eta(a_2^{-1}\alpha_2)+\eta(a_1^{-1}a_2^{-1}\alpha_1\alpha_2)\rt)\\ 
    & = q+ \eta(a_1)\sum_{\alpha_1\in\bbF_q}\eta(\alpha_1)+\eta(a_2)\sum_{\alpha_2\in\bbF_q}\eta(\alpha_2)+\eta(a_1a_2)\sum_{\substack{\alpha_1,\alpha_2\in\bbF_q\\ \alpha_1+\alpha_2=\alpha}}\eta(\alpha_1\alpha_2)\\ 
    & = q+\eta(a_1a_2)\sum_{\substack{\alpha_1,\alpha_2\in\bbF_q\\ \alpha_1+\alpha_2=\alpha}}\eta(\alpha_1\alpha_2) \by{\MultCharPropsthm[(i)]}\\ 
    & = q+\eta(a_1a_2)\sum_{\gm\in\bbF_q}\eta(\gm(\alpha-\gm))
  \end{align*}
  Now consider the polynomial $g(X)=\alpha X-X^2$. Then by \EtaQuadPolythm using the function $\vth$ we have $\sum_{\gm\in\bbF_q}\eta(g(\gm))=\vth(\alpha)\eta(-1)$. Therefore we have the lemma. 
\end{proof}
\begin{Theorem}{Solution Count for Non-degenerate Diagonal Quadratic Forms}{thm:qf-count-even-n-odd-q}
  \index{quadratic form!solution count}\index{solution count formula}%
  Let $f$ be a non-degenerate quadratic form in $\bbF_q[X_1,\dots,X_n]$ with $n$ even and $q$ odd. Then for any $\alpha\in\bbF_q$,
  \begin{enumerate}[label=(\roman*)]
    \item If $n$ is even $Z\lt(f(X_1,\dots,X_n)=\alpha\rt)=q^{n-1}+\vth(\alpha)\cdot q^{\nicefrac{(n-2)}{2}}\cdot \eta\lt((-1)^{\nicefrac{n}{2}}\cdot \det(f)\rt).$
    \item if $n$ is odd $Z\lt(f(X_1,\dots, X_n)=\alpha\rt)=q^{n-1}+q^{\nicefrac{(n-1)}{2}}\cdot \eta\lt((-1)^{\nicefrac{(n-1)}{2}}\cdot \alpha\cdot \det(f)\rt).$
  \end{enumerate}
\end{Theorem}
\begin{proof}
  By \QFDiagthm $f$ is equivalent to a diagonal quadratic form $g(X_1,\dots, X_n)=a_1X_1^2+\cdots a_nX_n^2$ by the non-singular linear transformation $M\in\bbF_q^{n\times n}$. By \QFEquivDetObs, we have $\eta(\det(f))=\eta(\det(g))$. Hence it suffices to prove the theorem for $g(X_1,\dots, X_n)=\alpha$. Since $f$ is non-degenerate by \QFDiagRankDetObs, $a_i\neq 0$ for all $i\in[n]$. 
  \begin{enumerate}[label=(\roman*)]
    \item Let $n$ is even. So take $m=\nicefrac{n}2$. Then 
  \begin{align*}
    Z(g(X_1,\dots, X_n)=\alpha) & = \sum_{\substack{\alpha_1,\dots,\alpha_m\in\bbF_q\\ \alpha_1+\cdots+\alpha_m=\alpha}}\prod_{j=1}^mZ\lt(a_{2j-1}X_{2j-1}^2+a_{2j}X_{2j}^2=\alpha_j\rt)\\ 
    & = \sum_{\substack{\alpha_1,\dots,\alpha_m\in\bbF_q\\ \alpha_1+\cdots+\alpha_m=\alpha}} \prod_{j=1}^m \lt(q+\vth(\alpha_j)\cdot\eta(-a_{2j-1}a_{2j})\rt)\by{\QFTwoVarlem}\\ 
    & = \sum_{\substack{\alpha_1,\dots,\alpha_m\in\bbF_q\\ \alpha_1+\cdots+\alpha_m=\alpha}} q^m+ \sum_{\substack{\alpha_1,\dots,\alpha_m\in\bbF_q\\ \alpha_1+\cdots+\alpha_m=\alpha}}\prod_{j=1}^m \vth(\alpha_j)\cdot\eta(-a_{2j-1}a_{2j})\\ 
    & = q^{m-1}\cdot q^m+ \eta\lt((-1)^m\det(g)\rt)\sum_{\substack{\alpha_1,\dots,\alpha_m\in\bbF_q\\ \alpha_1+\cdots+\alpha_m=\alpha}}\prod_{j=1}^m\vth(\alpha_j)\\ 
    & = q^{2m-1}+\eta\lt((-1)^m\det(g)\rt)\cdot \vth(\alpha)\cdot q^{m-1}\by{\QFConvProplem}\\ 
    & = q^{n-1}+\eta\lt((-1)^{\nicefrac{n}2}\det(f)\rt)\cdot \vth(\alpha)\cdot q^{\nicefrac{(n-2)}{2}}
  \end{align*}
  \item Now suppose $n$ is odd. We will use induction on $n$. Now this formula is valid for $n=1$. So suppose $n\geq 3$. Then 
  \begin{align*}
    Z(g(X_1,\dots, X_n)=\alpha) & = \sum_{\substack{\alpha_1,\alpha_2\in\bbF_q\\ \alpha_1+\alpha_2=\alpha}}Z(a_1X_1^2=\alpha_1)\cdot Z(a_2X_2^2+\cdots +a_nX_n^2=\alpha_2)\\ 
    & = \sum_{\substack{\alpha_1,\alpha_2\in\bbF_q\\ \alpha_1+\alpha_2=\alpha}} \lt(1+\eta(a_1^{-1}\alpha_1)\rt)\cdot \lt(q^{n-2}+\eta\lt((-1)^{\nicefrac{n-1}2}a_2\cdots a_n\rt)\cdot \vth(\alpha)\cdot q^{\nicefrac{(n-3)}{2}}\rt)\by{first part}\\ 
    & = \sum_{\substack{\alpha_1,\alpha_2\in\bbF_q\\ \alpha_1+\alpha_2=\alpha}}q^{n-1}+q^{n-2}\cdot \eta(a_1)\sum_{\alpha_1\in\bbF_q}\eta(\alpha_1)\\ 
    & \qquad\qquad\qquad\qquad\qquad\qquad+q^{\nicefrac{(q-3)}{2}}\cdot \eta\lt((-1)^{\nicefrac{n-1}2}a_2\cdots a_n\rt)\sum_{\alpha_2\in\bbF_q}\vth(\alpha_2)\\ 
    &\qquad\qquad\qquad\quad\ \; + q^{\nicefrac{(q-3)}{2}}\cdot \eta\lt((-1)^{\nicefrac{n-1}2}a_2\cdots a_n\rt)\sum_{\substack{\alpha_1,\alpha_2\in\bbF_q\\ \alpha_1+\alpha_2=\alpha}}\vth(\alpha_2)\eta(a_1^{-1}\alpha_1)\\ 
    & = q^{n-1}+ q^{\nicefrac{(q-3)}{2}}\cdot \eta\lt((-1)^{\nicefrac{n-1}2}a_1\cdots a_n\rt)\sum_{\substack{\alpha_1,\alpha_2\in\bbF_q\\ \alpha_1+\alpha_2=\alpha}}\vth(\alpha_2)\eta(\alpha_1)\by{\QFConvProplem}\\ 
    & = q^{n-1}+ q^{\nicefrac{(q-3)}{2}}\cdot \eta\lt((-1)^{\nicefrac{n-1}2}\cdot \det(f)\rt)\sum_{\gm\in\bbF_q}\eta(\gm)\cdot \vth(\alpha-\gm)
  \end{align*}So now we only have to calculate $\sum_{\gm\in\bbF_q}\eta(\gm)\cdot \vth(\alpha-\gm)$. Now $$\sum_{\gm\in\bbF_q}\eta(\gm)\cdot \vth(\alpha-\gm)=\sum_{\gm\in\bbF_q}\eta(\gm)\cdot [\vth(\alpha-\gm)+1]=q\cdot \eta(\alpha)$$
  \end{enumerate}
  So we have the results for both cases. 
\end{proof}

The following alternate proof uses Jacobi sums\index{Jacobi sum!connection to quadratic forms} directly and illustrates how the machinery of the previous chapter applies to count solutions of polynomial equations.


\begin{alternate-proof}
  Again it suffices to prove this theorem for the diagonal quadratic form $g(X_1,\dots, X_n)=a_1X_1^2+\cdots a_nX_n^2$ where $a_i\in\bbF_q^*$ for all $i\in[n]$. Here we will work with the extended definition of $\eta$. Let $\lm_0$ is the trivial multiplicative character and $\lm_1=\eta$. Then 
    \begin{align*}
      Z(g(X_1,\dots, X_n)=\alpha) & = \sum_{\substack{\alpha_1,\dots,\alpha_n\in\bbF_q\\ \alpha_1+\cdots+\alpha_n=\alpha}}\prod_{j=1}^nZ\lt(a_jX_j^2=\alpha_j\rt)\\ 
      & = \sum_{\substack{\alpha_1,\dots,\alpha_n\in\bbF_q\\ \alpha_1+\cdots+\alpha_n=\alpha}}\lt(\lm_1(a_j\alpha_j)+\lm_0(a_j\alpha_j)\rt)\\ 
      & = \sum_{\substack{\alpha_1,\dots,\alpha_n\in\bbF_q\\ \alpha_1+\cdots+\alpha_n=\alpha}}\sum_{i_1,\dots,i_n=0}^1\lm_{i_1}(a_1\alpha_1)\cdots \lm_{i_n}(a_n\alpha_n)\\ 
      & = \sum_{i_1,\dots,i_n=0}^1 \lt(\prod_{j=1}^n\lm_{i_j}(a_j)\rt) \sum_{\substack{\alpha_1,\dots,\alpha_n\in\bbF_q\\ \alpha_1+\cdots+\alpha_n=\alpha}}\prod_{j=1}^n\lm_{i_j}(\alpha_j)\\ 
      & = \sum_{i_1,\dots,i_n=0}^1 \lt(\prod_{j=1}^n\lm_{i_j}(a_j)\rt) \cdot J_{\alpha}\lt(\lm_{i_1},\dots, \lm_{i_n}\rt)
    \end{align*}
    Now if not all but some $\lm_{i_j}$'s are trivial then $J_{\alpha}\lt(\lm_{i_1},\dots, \lm_{i_n}\rt)=0$ by \JacobiCompletethm[(ii)]. So this remains the cases when all are $\lm_{i,j}$'s are trivial or non-trivial. Now if $i_j=0$ for all $j\in[n]$ then by \JacobiCompletethm[(i)], $J_{\alpha}\lt(\lm_{i_1},\dots, \lm_{i_n}\rt)=q^{n-1}$. Therefore by \begin{equation}\label{eq:quad-jacobi-etas}
      Z(g(X_1,\dots, X_n)=\alpha) =q^{n-1}+\eta(\det(f))\cdot J_{\alpha}\underbrace{(\eta,\dots, \eta)}_{n-\text{times}}
    \end{equation}If $\alpha\neq 0$ then above expression gives us \begin{equation}\label{eq:quad-jacobi-nonzero}
      Z(g(X_1,\dots, X_n)=\alpha) =q^{n-1}+\eta(\det(f))\cdot \eta^n(\alpha)\cdot J(\eta,\dots, \eta)
    \end{equation}
    \begin{enumerate}[label=(\roman*)]
      \item Now suppose $n$ is even. Let $\alpha\neq 0$. Since $n$ is even $\eta^n$ is trivial character. So from the equation  \eqref{eq:quad-jacobi-nonzero} we have $$Z(g(X_1,\dots, X_n)=\alpha) =q^{n-1}+\eta(\det(f))\cdot J(\eta,\dots, \eta)$$ Let $\chi\in\Xq$ is a non-trivial additive character of $\bbF_q$. Then by \JGrelthm[(ii)] we have $$J(\eta,\dots, \eta) = -\frac1qG(\eta,\chi)^n = -\frac1q\lt(G^2(\eta,\chi)\rt)^{\nicefrac{n}{2}} = -\frac1q(\eta(-1)q)^{\nicefrac{n}2} = -q^{\nicefrac{(n-2)}{2}}\cdot \eta\lt((-1)^{\nicefrac{n}2}\rt)$$Hence if $\alpha\neq 0$ then we have $$Z(g(X_1,\dots, X_n)=\alpha)=q^{n-1}+\eta(\det(f))\cdot \lt(-q^{\nicefrac{(n-2)}{2}}\cdot \eta\lt((-1)^{\nicefrac{n}2}\rt)\rt)=q^{n-1}-q^{\nicefrac{(q-2)}{2}}\eta\lt((-1)^{\nicefrac{q}{2}}\cdot \det(f)\rt)$$
      Now suppose $\alpha= 0$. Then \eqref{eq:quad-jacobi-etas} gives $$Z(g(X_1,\dots, X_n)=\alpha) =q^{n-1}+\eta(\det(f))\cdot J_{0}(\eta,\dots, \eta)$$Since $n$ is even $\eta^n$ is trivial. Therefore by \JacZeroRellem[(ii)] we have $$J_0\underbrace{(\eta,\dots, \eta)}_{n-\text{times}}=\eta(-1)\cdot (q-1)\cdot J\underbrace{(\eta,\dots\eta)}_{(n-1)-\text{times}}$$So by \JGrelthm[(i)] we have $$J\underbrace{(\eta,\dots\eta)}_{(n-1)-\text{times}}=\frac{G(\eta,\chi)^{n-1}}{G(\eta^{n-1},\chi)}=G^{n-2}(\eta,\chi)=\lt(G^2(\eta,\chi)\rt)^{\nicefrac{(n-2)}{2}}= q^{\nicefrac{(n-2)}{2}}\cdot \eta\lt((-1)^{\nicefrac{(n-2)}{2}}\rt)$$Therefore from \eqref{eq:quad-jacobi-etas} we get $$Z(g(X_1,\dots, X_n)=\alpha)=q^{n-1}+(q-1)\cdot q^{\nicefrac{(n-2)}{2}}\cdot \eta\lt((-1)^{\nicefrac{n}{2}}\cdot \det(f)\rt)$$
      So combining the result of $\alpha\neq 0$ and $\alpha=0$ we get the theorem. 
      \item Now suppose $n$ is odd. Then $\eta^n=\eta$ which is non-trivial. If $\alpha=0$ then by \JacZeroRellem[(i)] we have $J_0(\eta,\dots,\eta)=0$. So assume $\alpha\neq 0$. Then in the equation \eqref{eq:quad-jacobi-nonzero} by using \JGrelthm[(i)] we have
      $$J(\eta,\dots,\eta)=\frac{G(\eta,\chi)^{n}}{G(\eta^{n},\chi)}=G^{n-1}(\eta,\chi)=\lt(G^2(\eta,\chi)\rt)^{\nicefrac{(n-1)}{2}}= q^{\nicefrac{(n-1)}{2}}\cdot \eta\lt((-1)^{\nicefrac{(n-1)}{2}}\rt)$$So in both cases $$Z(g(X_1,\dots, X_n)=\alpha) =q^{n-1}+ q^{\nicefrac{(n-1)}{2}}\cdot \eta\lt((-1)^{\nicefrac{(n-1)}{2}}\cdot \alpha\cdot \det(f)\rt)$$This expression works for $\alpha=0$ because then the second term becomes zero. 
    \end{enumerate}
    Hence we have theorem. 
\end{alternate-proof}

\begin{note}
  For any general degree-$2$ polynomial $f\in\bbF_q[X_1,\dots,X_n]$ there exists $g,h\in\bbF_q[X_1,\dots,X_n]$ where $g$ is a quadratic form and $\deg(h)\leq 1$. Then by some non-singular linear transformation we can transform $f$ into an equivalent diagonal quadratic form. Then we obtain an equation of following form:
  $$\tdf(X_1,\dots, X_n)=a_1X_1^2+\cdots a_kX_k^2+b_1X_1+\cdots b_nX_n$$ where $a_i\in\bbF_q^*$ for all $i\in[k]$, $k\leq n$ and $b_j\in\bbF_q$ for all $j\in[n]$. Hence for any $\alpha\in\bbF_q$, the number of solutions of $f(X_1,\dots, X_n)=\alpha$ is same as number of solutions of $\tdf(X_1,\dots, X_n)=\alpha$. 

  If $k<n$ then without loss of generality we can assume $b_n\neq 0$. Then the number of solutions is $q^{n-1}$ as we can arbitrarily substitute elements for $X_1,\dots, X_{n-1}$ and then value of $X_n$ is uniquely determined. 
  
  If $k=n$ then takes the non-singular linear transformation $X_i=Y_i-b_i(2a_i)^{-1}$ for all $i\in[n]$ gives an equivalent diagonal quadratic form $a_1Y_1^2+\cdots +a_nY_n^2=c$ for some $c\in\bbF_q$. Now we can use \QFCountOddthm to find the number of solutions. 
\end{note}

\subsection{Even Characteristic Quadratic Forms}
\index{quadratic form!even characteristic}%
We will now study quadratic forms and its number of solutions over finite fields $\bbF_q$ with $\Char(\bbF_q)=2$ i.e. $q$ is even. Again like in the case of odd characteristic we will first reduce any general quadratic form to canonical forms and then we will try to find the number of solutions of such canonical forms. 

For even characteristic fields we define \emph{non-degenerate} quadratic forms differently. Let $f\in\bbF_q[X_1,\dots,X_n]$ is a quadratic form in $n$-variables. Then $f$ is called non-degenerate if $f$ is not equivalent to a quadratic form in fewer than $n$ variables.
\begin{note}
  This definition is consistent with the odd-characteristic case: there, if $f$ is equivalent to a form in fewer than $n$ variables, \QFDiagthm and \QFDiagRankDetObs show that $\Rank(C_f)<n$, i.e.\ $\det(f)=0$, which is the standard criterion for degeneracy. So in both characteristics one can reduce to the non-degenerate case by passing to the equivalent form in fewer variables.
\end{note}  
\begin{lemma}{Reduction for Quadratic Forms in even characteristic}{lem:qf-reduction-even-char}
  A non-degenerate quadratic form $f\in\bbF_q[X_1,\dots,X_n]$ with $q$ even and $n\geq 3$ is equivalent to $X_1X_2+g(X_3,\dots, X_n)$ where $g$ is a quadratic form in $n-2$ variables.
\end{lemma}
\begin{proof}
  First we will show that $f$ is equivalent to a quadratic form where coefficient of $X_1^2$ is $0$. So let $$f(X_1,\dots, X_n)=\sum_{i,j\in[n]}a_{i,j}X_iX_j$$ where $a_{i,j}=a_{j,i}\in\bbF_q$. If there exists $i\in[n]$ such that $a_{i,i}=0$ then we can permute the variables so that $X_i$ goes to $X1$ and thus $a_{1,1}=0$. So assume $a_{i,i}\neq 0$ for all $i\in[n]$. Now if $a_{i,j}=0$ for all $i\neq j$ then $$f(X_1,\dots, X_n)=a_{1,1}X_1^2+\cdots +a_{n,n}X_n^2=\lt(a_{1,1}^{\nicefrac{q}2}X_1+\cdots +a_{n,n}^{\nicefrac{q}2}X_n\rt)^2$$then $f$ is equivalent to $Y_1^2$ by the non-singular linear transformation\begin{align*}
    Y_1&=a_{1,1}^{\nicefrac{q}2}X_1+\cdots +a_{n,n}^{\nicefrac{q}2}X_n&&\\
    Y_i&=X_i\ \forall\ i\in\{2,\dots, n\}&&
  \end{align*}This contradicts the non-degenerate property of $f$ \ctr So there exists $i,j\in[n]$, $i\neq j$ such that $a_{i,j}\neq 0$. Without loss of generality suppose $a_{2,3}\neq 0$. Now $$f(X_1,\dots, X_n)=a_{2,2}X_2+X_2(a_{1,2}X_1+a_{2,3}X_3+\cdots +a_{2,n}X_n)+g_1(X_1,X_3\dots, X_n)$$Then take the non-singular linear transformation
  \begin{align*}
    Y_3&=a_{2,3}^{-1}(a_{1,2}X_1+a_{2,3}X_3+\cdots +a_{2,n}X_n)&&\\
    Y_i&=X_i\ \forall\ i\in[n]\setminus\{3\}&&
  \end{align*}This gives an equivalent quadratic form $a_{22}Y_2^2+Y_2Y_3+g_2(Y_1,Y_3,\dots, Y_n)$. Let $b$ is the coefficient of $X_1^2$ in $g_2$. So now take the non-singular linear transformation
  \begin{align*}
    Z_2&=\lt(a_{2,2}^{-1}b\rt)^{\nicefrac{q}2}Y_1+Y_2&&\\
    Z_i&=Y_i\ \forall\ i\in[n]\setminus\{2\}&&
  \end{align*}This transformation makes the coefficient of $Z_1^2$ to be $0$. Now we'll use $f$ to denote this final quadratic form with variables $X_1,\dots, X_n$. Since $f$ is non-degenerate not all $a_{i,j}$'s are $0$ where $i\neq j$. Without loss of generality suppose $a_{1,2}\neq 0$. Then take the non-singular linear transformation
    \begin{align*}
    Y_2&=a_{1,2}^{-1}(Y_2+a_{1,3}Y_3+\cdots +a_{1,n}Y_n)&&\\
    Y_i&=X_i\ \forall\ i\in[n]\setminus\{2\}&&
  \end{align*}Then we get a quadratic form of the type $$Y_1Y_2+\sum_{2\leq i,j\leq n}c_{i,j}Y_iY_j$$ where $c_{i,j}\in\bbF_q$ for all $2\leq i,j\leq n$. So take the non-singular linear transformation   \begin{align*}
    Z_1& = Y_1+c_{2,2}Y_2+\cdots +c_{2,n}Y_n&&\\
    Z_i&=Y_i\ \forall\ i\in[n]\setminus\{1\}&&
  \end{align*}
  This transformation gives a quadratic form of the type $Z_1Z_2+g_3(Z_3,\dots, Z_n)$. So we have the lemma. 
\end{proof}

So now we will use this theorem to find proper canonical forms for any general quadratic forms. But we have to treat the cases of odd and even number of variables separately. 
\begin{Theorem}{Canonical Forms for even characteristic}{thm:qf-canonical-even-char}
  Let $f\in\bbF_q[X_1,\dots,X_n]$ be a non-degenerate quadratic form with $q$ even.
  \begin{enumerate}[label=(\roman*)]
    \item If $n$ is odd then $f$ is equivalent to $X_1X_2+X_3X_4+\cdots +X_{n-2}X_{n-1}+X_n^2$.
    \item If $n$ is even then $f$ is \begin{enumerate}
      \item either equivalent to $X_1X_2+\cdots+X_{n-1}X_n$, 
      \item or to $X_1X_2+\cdots+X_{n-3}X_{n-2}+X_{n-1}^2+a\cdot X_n^2$ where $a\in\bbF_q$ satisfies $\tr(a)=1$ where $\tr$ is the absolute trace function from $\bbF_q$ to $\bbF_2$. 
    \end{enumerate}
  \end{enumerate}
\end{Theorem}
\begin{proof}
  \begin{enumerate}[label=(\roman*)]
    \item Suppose $n$ is odd. By applying \QFRedEvenlem again and again on $f$ we get an equivalent quadratic form of the type $X_1X_2+X_3X_4+\cdots +X_{n-2}X_{n-1}+aX_n^2$ for some $a\in\bbF_q^*$. Then take the non-singular linear transformation
    \begin{align*}
      Y_n& = a^{-\nicefrac{q}{2}}X_n&&\\ 
      Y_i& = X_i\ \forall\ i\in[n-1]&&
    \end{align*}
    Then we get the desired quadratic form. 
    \item So now assume $n$ is even. Again by applying \QFRedEvenlem again and again on $f$ we get an equivalent quadratic form of the type $$X_1X_2+\cdots+X_{n-3}X_{n-2}+bX_{n-1}^2+cX_{n-1}X_n+dX_n^2$$with $b,c,d\in\bbF_q$. If $c=0$ then take the non-singular linear transformation
    \begin{align*}
      Y_{n-1}& = b^{\nicefrac{q}2}X_{n-1}+d^{\nicefrac{q}2}X_n&&\\ 
      Y_i& = X_i\ \forall\ i\in[n]\setminus\{n-1\}&&
    \end{align*}Then we get a quadratic form with $(n-1)$-variables. But $f$ is non-degenerate. Hence contradiction \ctr Therefore $c\in\bbF_q^*$. 
    
    If $b=0$ then take the non-singular linear transformation 
      \begin{align*}
      Y_{n-1}& = cX_{n-1}+dX_n&&\\ 
      Y_i& = X_i\ \forall\ i\in[n]\setminus\{n-1\}&&
    \end{align*}With this transformation we get the first form. 
    
    
    If $b\neq 0$ then take the non-singular linear transformation
    \begin{align*}
      Y_{n-1}& = b^{-\nicefrac{q}2}X_{n-1}&&\\
      Y_n & = b^{\nicefrac{q}2}c^{-1}X_n && \\ 
      Y_i& = X_i\ \forall\ i\in[n-2]&&
    \end{align*}Then we get a quadratic form of the type $$Y_1Y_2+\cdots+Y_{n-3}Y_{n-2}+Y_{n-1}^2+Y_{n-1}Y_n+aY_n^2$$for some $a\in\bbF_q$. If $X^2+X+a$ is reducible in $\bbF_q$ then there exists $c_1,c_2\in\bbF_q$ such that $$X^2+X+a=(X+c_1)(X+c_2)\implies Y_{n-1}^2+Y_{n-1}Y_n+aY_n^2=(Y_{n-1}+c_1Y_n)(Y_{n-1}+c_2Y_n)$$Therefore if we take the non-singular linear transformation 
    \begin{align*}
      Z_{n-1}& = Y_{n-1}+c_1Y_n&&\\
      Z_n & = Y_{n-1}+c_2Y_n && \\ 
      Z_i& = Y_i\ \forall\ i\in[n-2]&&
    \end{align*}then we get a quadratic form of the first form. Now $X^2+X+a$ is irreducible in $\bbF_q[X]$ if  $\tr(a)=1$. Then we have the second form. 
  \end{enumerate}
  Therefore we have the canonical forms for all cases. 
\end{proof}
Now by the above theorem it suffices to only find the number of solutions for canonical non-degenerate quadratic forms. We will first show two results on two variable version then we will use it to show the general version. 
\begin{lemma}{Solutions of a Two-Variable Quadratic Form in even characteristic}{lem:qf-two-var-count-even}
  For even $q$, let $a,\alpha\in\bbF_q$ with $\tr(a)=1$. Then $$Z\lt(X_1^2+X_1X_2+aX_2^2=\alpha\rt)=q-\vth(\alpha).$$where $\tr$ is the absolute trace function from $\bbF_q$ to $\bbF_2$. 
\end{lemma}
\begin{proof}
  Since $\tr(a)=1$ the polynomial $X^2+X+a$ is irrducible in $\bbF_q[X]$. So there exists $\gm\in\bbF_{q^2}\setminus\bbF_q$ such that $X^2+X+a=(X+\gm)(X+\gm^q)$. So $$f(X_1,X_2)=X_1^2+X_1X_2+aX_2^2=(X_1+\gm X_2)(X_1+\gm^q X_2)$$Therefore for any $(\beta_1,\beta_2)\in\bbF_q^2$ we have $$f(\beta_1,\beta_2)=(\beta_1+\gm\beta_2)(\beta_1+\gm^q\beta_2)=(\beta_1+\gm\beta_2)(\beta_1+\gm\beta_2)^q=(\beta_1+\gm\beta_2)^{q+1}$$So we get $$Z(f(X_1,X_2)=\alpha)=|\{\gm\in\bbF_{q^2}\colon \gm^{q+1}=\alpha\}|$$Hence if $\alpha=0$ then $Z\lt(X_1^2+X_1X_2+aX_2^2=0\rt)=1=q-\vth(0)$. If $\alpha\neq 0$ since $\bbF_{q^2}^*$ is cyclic we have $\alpha^{\nicefrac{(q^2-1)}{q+1}}=\alpha^{q-1}=1$ and hence there are $q+1$ elements in $\bbF_{q^2}$ such that $\gm^{q+1}=\alpha$. Hence $Z\lt(X_1^2+X_1X_2+aX_2^2=\alpha\rt)=q+1=q+\vth(\alpha)$.
\end{proof}
\begin{lemma}{}{lem:x1x2-quad-form}
  For even $q$, let $\alpha\in\bbF_q$. Then $Z(X_1X_2=\alpha)=q+\vth(\alpha)$. 
\end{lemma}
\begin{proof}
  If $\alpha\neq 0$ then both $X_1$ and $X_2$ has to be assigned with non-zero values. After assigning $X_1$ arbitrarily from $\bbF_q^*$ the value of $X_2$ gets uniquely determined. Therefore $Z(X_1X_2=\alpha)=q-1=q+\vth(\alpha)$. 

  Suppose $\alpha=0$. Then at least one of $X_1$ or $X_2$ gets assigned $0$. Then the other variable can take any value from $\bbF_q^*$. Therefore there are $2q-1$ such solutions. Hence $Z(X_1X_2=\alpha)=2q-1=q+\vth(\alpha)$. So we have the lemma. 
\end{proof}
\begin{Theorem}{Solution Counts for Canonical Quadratic Forms in even characteristic}{thm:qf-count-even-char}
  Let $\bbF_q$ be a finite field with $q$ even. Then for any $\alpha\in\bbF_q$:
  \begin{enumerate}[label=(\roman*)]
    \item If $n$ is odd, $Z\lt(X_1X_2+\cdots + X_{n-2}X_{n-1}+X_n^2=\alpha\rt)=q^{n-1}+\vth(\alpha)\cdot q^{\nicefrac{(n-1)}{2}}.$
    \item If $n$ is even:
    \begin{enumerate}[label=(\alph*)]
      \item $Z\lt(X_1X_2+\cdots +X_{n-1}X_n=\alpha\rt)=q^{n-1}+\vth(\alpha)\cdot q^{\nicefrac{(n-2)}{2}}$;
      \item for $a\in\bbF_q$ with $\tr(a)=1$, then $Z\lt(X_1X_2+\cdots+X_{n-1}X_{n}+X_{n-1}^2+a\cdot X_n^2=\alpha\rt)=q^{n-1}-\vth(\alpha)\cdot q^{\nicefrac{(n-2)}{2}}$, where $\tr$ is the absolute trace function from $\bbF_q$ to $\bbF_2$.
    \end{enumerate}
  \end{enumerate}
\end{Theorem}
\begin{proof}
  \begin{enumerate}[label=(\roman*)]
    \item Suppose $n$ is odd. Since $q$ is even the equation $X^2=\alpha$ for any $\alpha\in\bbF_q$ has an unique solution in $\bbF_q$. Therefore $Z\lt(X_1X_2+\cdots + X_{n-2}X_{n-1}+X_n^2=\alpha\rt)=q^{n-1}+\vth(\alpha)\cdot q^{\nicefrac{(n-1)}{2}}$ because after assigning the variables $X_1,\dots, X_{n-1}$ arbitrarily the value of $X_n$ is uniquely determined. 
    \item So let $n$ is even. Let $m=\nicefrac{n}2$. Then we have \begin{align*}
        Z\lt(X_1X_2+\cdots +X_{n-1}X_n=\alpha\rt)& = \sum_{\substack{\alpha_1,\dots,\alpha_m\in\bbF_q\\ \alpha_1+\cdots+\alpha_m=\alpha}}\prod_{j=1}^n Z(X_{2j-1}X_{2j}=\alpha_j)\\ 
        & = \sum_{\substack{\alpha_1,\dots,\alpha_m\in\bbF_q\\ \alpha_1+\cdots+\alpha_m=\alpha}} \prod_{j=1}^m (q+\vth(\alpha_j)) \byl{x1x2-quad-form}\\ 
        & = q^{2m-1}+\sum_{\substack{\alpha_1,\dots,\alpha_m\in\bbF_q\\ \alpha_1+\cdots+\alpha_m=\alpha}}\prod_{j=1}^m\vth(\alpha_j)\\ 
        & = q^{2m-1}+\vth(\alpha)\cdot q^{m-1}=q^{n-1}+\vth(\alpha)\cdot q^{\nicefrac{(n-2)}{2}}\by{\QFConvProplem}
      \end{align*}
      So we have the first part. Now for the second part we already showed the $n=2$ case in \QFTwoVarEvenlem. So assume $n\geq 4$.  
      \begin{align*}
        & Z\lt(X_1X_2+\cdots+X_{n-3}X_{n-2}+X_{n-1}^2+a\cdot X_n^2=\alpha\rt)\\ 
        = \;&\;\sum_{\substack{\alpha_1,\alpha_2\in\bbF_q\\ \alpha_1+\alpha_2=\alpha}}Z\lt(X_1X_2+\cdots+X_{n-3}X_{n-2}=\alpha_1\rt)\cdot Z(X_{n-1}X_{n}+X_{n-1}^2+a\cdot X_n^2=\alpha_2)\\ 
        = \;&\; \sum_{\substack{\alpha_1,\alpha_2\in\bbF_q\\ \alpha_1+\alpha_2=\alpha}} \lt(q^{n-3}+\vth(\alpha_1)\cdot q^{\nicefrac{(n-4)}{2}}\rt)(q-\vth(\alpha_2))\by{first part and \QFTwoVarEvenlem}\\ 
        = \;&\; q^{n-1}-q^{\nicefrac{(n-4)}{2}}\sum_{\substack{\alpha_1,\alpha_2\in\bbF_q\\ \alpha_1+\alpha_2=\alpha}}\vth(\alpha_1)\cdot\vth(\alpha_2) \by{\QFConvProplem}\\ 
        = \;&\; q^{n-1}-q^{\nicefrac{(n-4)}{2}}\cdot q\cdot \vth(\alpha)\\ 
        = \;&\; q^{n-1}-\vth(\alpha)\cdot q^{\nicefrac{(n-2)}{2}}
      \end{align*}
      So now we have the result for second part
  \end{enumerate}
  Therefore we get the values of number of solutions of quadratic forms in even characteristic finite fields. 
\end{proof}
\section{Diagonal Equations}\label{sec:diagonal-equations}
In the previous section we talked about equations involving diagonal quadratic forms. These are a special case of much general family of equations called \emph{diagonal equations}
\begin{Definition}{Diagonall Equations}{}
  A diagonal equation over $\bbF_q$ is an equation of the type $$a_1X_1^{k_1}+\cdots+a_nX_n^{k_n}=\alpha$$ for some positive integers $k_1,\dots,k_n\in\bbN$ with coefficients $a_1,\dots,a_n\in\bbF_q^*$ and $\alpha\in\bbF_q$. 
\end{Definition}


\subsection{Counting Solutions using Jacobi Sums}
Now we express the number of solutions $Z\lt(a_1X_1^{k_1}+\cdots+a_nX_n^{k_n}=\alpha\rt)$ in terms of Jacobi sums. 
$$Z\lt(a_1X_1^{k_1}+\cdots+a_nX_n^{k_n}=\alpha\rt) = \sum_{\substack{\alpha_1,\dots,\alpha_n\in\bbF_q\\ \alpha_1+\cdots +\alpha_n=\alpha}}\prod_{j=1}^n Z\lt( a_jX_j^{k_j}=\alpha_j\rt)=\sum_{\substack{\alpha_1,\dots,\alpha_n\in\bbF_q\\ \alpha_1+\cdots +\alpha_n=\alpha}}\prod_{j=1}^n Z\lt(X_j^{k_j}=a_j^{-1}\alpha_j\rt)$$ Let $\lm$ is a multiplicative character of $\bbF_q$ with $\ord(\lm)=d=\gcd(k,q-1)$. Then for any $k\in\bbN$ and $\gm\in\bbF_q$ we have $Z(X^k=\gm)=\sum_{j=0}^{d-1}\lm^j$. So here let for $j\in[n]$ let $d_i=\gcd(k_i,q-1)$ and $\lm_i$ is the multiplicative character of $\bbF_q$ of order $d_i$. Then we have 
\begin{align*}
  Z\lt( a_jX_j^{k_j}=\alpha_j\rt) & = \sum_{\substack{\alpha_1,\dots,\alpha_n\in\bbF_q\\ \alpha_1+\cdots +\alpha_n=\alpha}}\prod_{j=1}^n Z\lt(X_j^{k_j}=a_j^{-1}\alpha_j\rt)\\ 
  & = \sum_{\substack{\alpha_1,\dots,\alpha_n\in\bbF_q\\ \alpha_1+\cdots +\alpha_n=\alpha}} \lt(\sum_{t_1=0}^{d_1-1}\lm_1^{t_1}(a_1^{-1}\alpha_1)\rt)\cdots\lt(\sum_{t_n=0}^{d_n-1}\lm_n^{t_n}(a_n^{-1}\alpha_n)\rt)\\
  & = \sum_{t_1=0}^{d_1-1}\cdots \sum_{t_n=0}^{d_n-1}\lm_1^{t_1}(a_1^{-1})\cdots\lm_n^{t_n}(a_n^{-1}) \sum_{\substack{\alpha_1,\dots,\alpha_n\in\bbF_q\\ \alpha_1+\cdots +\alpha_n=\alpha}}\lm_1^{t_1}(\alpha_1)\cdots\lm_n^{t_n}(\alpha_n)\\
  & = \sum_{t_1=0}^{d_1-1}\cdots \sum_{t_n=0}^{d_n-1}\ov{\lm}_1^{t_1}(a_1)\cdots\ov{\lm}_n^{t_n}(a_n) \cdot J_{\alpha}\lt(\lm_1^{t_1},\dots, \lm_n^{t_n}\rt)
\end{align*}
If $(t_1,\dots, t_n)=(0,\dots, o)$ then by \JacobiCompletethm[(i)] we have $J_{\alpha}(\lm_1^{t_1},\dots,\lm_n^{t_n})=q^{n-1}$. Also by \JacobiCompletethm[(ii)] if at least one of $t_j$ is zero then $J_{\alpha}(\lm_1^{t_1},\dots,\lm_n^{t_n})=0$. So we have $$Z\lt( a_jX_j^{k_j}=\alpha_j\rt)=q^{n-1}+\sum_{t_1=1}^{d_1-1}\cdots \sum_{t_n=1}^{d_n-1}\ov{\lm}_1^{t_1}(a_1)\cdots\ov{\lm}_n^{t_n}(a_n) \cdot J_{\alpha}\lt(\lm_1^{t_1},\dots, \lm_n^{t_n}\rt)$$So we have the following observation:
\begin{observation}\label{obs:diagonal-jacobi-formula}
  For any $\alpha\in\bbF_q$ let $a_1X_1^{k_1}+\cdots+a_nX_n^{k_n}=\alpha$ is a diagonal equation with $k_1,\dots,k_n\in\bbN$ and $a_1,\dots,a_n\in\bbF_q^*$. Suppose for all $i\in[n]$, $d_i=\gcd(k_i,q-1)$ and $\lm_i\in\Mq$ be a multiplicative character of $\bbF_q$ of order $d_i$, then \begin{equation}\label{eq:diagonal-solutions}
    Z\lt( a_jX_j^{k_j}=\alpha_j\rt)=q^{n-1}+\sum_{t_1=1}^{d_1-1}\cdots \sum_{t_n=1}^{d_n-1}\ov{\lm}_1^{t_1}(a_1)\cdots\ov{\lm}_n^{t_n}(a_n) \cdot J_{\alpha}\lt(\lm_1^{t_1},\dots, \lm_n^{t_n}\rt)
  \end{equation}
\end{observation}

Now we distinguish the cases when $\alpha=0$ and when $\alpha\neq 0$. If $\alpha=0$ then by \JacZeroRellem[(i)] for any $(t_1,\dots, t_n)$ if $\lm_1^{t_1}\cdots \lm_n^{t_n}$ is non-trivial then $J_0\lt(\lm_1^{t_1},\dots ,\lm_n^{t_n}\rt)=0$. So let $T$ denote the set of all tuples $(t_1,\dots, t_n)\in \bigtimes_{i=1}^n [d_i-1]$ such that $\lm_1^{t_1}\cdots \lm_n^{t_n}$ is trivial. Then we have the following result
\begin{Theorem}{Solution Count for Diagonal Equations at Zero}{thm:diagonal-count-zero}
  For any diagonal equation $a_1X_1^{k_1}+\cdots+a_nX_n^{k_n}=0$ with $k_1,\dots,k_n\in\bbN$ and $a_1,\dots,a_n\in\bbF_q^*$, the number of solutions is $$Z\lt(a_1X_1^{k_1}+\cdots+a_nX_n^{k_n}=0\rt)=q^{n-1}+\sum_{(t_1,\dots, t_n)\in T}\ov{\lm}_1^{t_1}(a_1)\cdots\ov{\lm}_n^{t_n}(a_n) \cdot J_0\lt(\lm_1^{t_1},\dots, \lm_n^{t_n}\rt),$$where $d_i=\gcd(k_i,q-1)$, $\lm_i$ has order $d_i$, and $T$ is the set of tuples $(t_1,\dots,t_n)$ with each $t_i\geq 1$ such that $\lm_1^{t_1}\cdots\lm_n^{t_n}$ is trivial.
\end{Theorem}
Now we can a trivial upper bound on the set $T$. Since $T\subseteq \bigtimes_{i=1}^n [d_i-1]$ we have $|T|\leq \prod_{i=1}^n (d_i-1)$. So we get the following upper bound from the above theorem.
\begin{corolary}{Bound on Solutions of Diagonal Equation at Zero}{cor:diagonal-bound-zero}
  For any diagonal equation $a_1X_1^{k_1}+\cdots+a_nX_n^{k_n}=0$ with $k_1,\dots,k_n\in\bbN$ and $a_1,\dots,a_n\in\bbF_q^*$, $$\lt|Z\lt(a_1X_1^{k_1}+\cdots+a_nX_n^{k_n}=0\rt)-q^{n-1}\rt|\leq (q-1)\cdot q^{\nicefrac{n}{2}-1}\prod_{i=1}^n(d_i-1),$$where $d_i=\gcd(k_i,q-1)$ for all $i\in[n]$.
\end{corolary}
\begin{proof}
  For all $(t_1,\dots, t_n)\in T$ we have $\lm_1^{t_1}\cdots \lm_n^{t_n}$ is trivial. Therefore $|J_0(\lm_1^{t_1},\dots, \lm_n^{t_n})|=(q-1)\cdot q^{\nicefrac{(n-2)}{2}}$ by \JacobiCompletethm[(iv)] Therefore using the trivial bound of $T$ and absolute value of $|J_0(\lm_1^{t_1},\dots, \lm_n^{t_n})|$ we get the above result. 
\end{proof}

Let $\lm$ is the generator of $\Mq$. Then for all $i\in[n]$, $\lm_i=\lm^{\nicefrac{(q-1)}{d_i}}$. Hence $$\lm_1^{t_1}\cdots \lm_n^{t_n}=\lm^{t_1\frac{q-1}{d_1}+\cdots+ t_n\frac{q-1}{d_n}}=\lm^{(q-1)\lt(\frac{t_1}{d_1}\cdots+\frac{t_n}{d_n}\rt)}$$ Therefore $\lm_1^{t_1}\cdots \lm_n^{t_n}$ is trivial if and only if $\sum_{i=1}^n \frac{t_i}{d_i}\in\bbZ$. So let $M(d_1,\dots, d_n)$ is the number of all tuples $(t_1,\dots,t_n)\in\bigtimes_{i=1}^n [d_i-1]$ such that $\sum_{i=1}^n \frac{t_i}{d_i}\in\bbZ$. Therefore $|T|=M(d_1,\dots, d_n)$. So we have the following theorem. 
\begin{corolary}{Bound on Solutions of Diagonal Equation at Zero (via $M$)}{cor:diagonal-bound-zero-M}
  For any diagonal equation $a_1X_1^{k_1}+\cdots+a_nX_n^{k_n}=0$ with $k_1,\dots,k_n\in\bbN$ and $a_1,\dots,a_n\in\bbF_q^*$,
  $$\lt|Z\lt(a_1X_1^{k_1}+\cdots+a_nX_n^{k_n}=0\rt)-q^{n-1}\rt|\leq (q-1)\cdot q^{\nicefrac{(n-2)}{2}}\cdot M(d_1,\dots, d_n),$$where $d_i=\gcd(k_i,q-1)$ for all $i\in[n]$.
\end{corolary}

So now assume $\alpha\neq 0$ then we have $J_{\alpha}(\lm_1^{t_1},\dots, \lm_n^{t_n})=\lt(\lm_1^{t_1}\cdots\lm_n^{t_n}\rt)(\alpha)J(\lm_1^{t_1},\dots, \lm_n^{t_n})$. So we get the following formulation of number of solutions for $\alpha\neq 0$. 
\begin{Theorem}{Solution Count for Diagonal Equations at Nonzero Values}{thm:diagonal-count-nonzero}
  Let $\alpha\in\bbF_q^*$. For any diagonal equation $a_1X_1^{k_1}+\cdots+a_nX_n^{k_n}=\alpha$ with $k_1,\dots,k_n\in\bbN$ and $a_1,\dots,a_n\in\bbF_q^*$,
  $$Z\lt(a_1X_1^{k_1}+\cdots+a_nX_n^{k_n}=\alpha\rt)=q^{n-1}+\sum_{t_1=1}^{d_1-1}\cdots \sum_{t_n=1}^{d_n-1}\lm_1^{t_1}(a_1^{-1}\alpha)\cdots\lm_n^{t_n}(a_n^{-1}\alpha) \cdot J\lt(\lm_1^{t_1},\dots, \lm_n^{t_n}\rt),$$where $d_i=\gcd(k_i,q-1)$ and $\lm_i\in\Mq$ has order $d_i$ for all $i\in[n]$.
\end{Theorem}

Now we will show another way to find the number of solutions for $\alpha\neq 0$ using the notation $M(d_1,\dots, d_n)$ we used in the case of $\alpha=0$. 
\begin{Theorem}{Bound on Solutions of Diagonal Equation at Nonzero Values}{thm:diagonal-bound-nonzero}
  Let $\alpha\in\bbF_q^*$. For any diagonal equation $a_1X_1^{k_1}+\cdots+a_nX_n^{k_n}=\alpha$ with $k_1,\dots,k_n\in\bbN$ and $a_1,\dots,a_n\in\bbF_q^*$,
  $$\lt|Z\lt(a_1X_1^{k_1}+\cdots+a_nX_n^{k_n}=\alpha\rt)-q^{n-1}\rt| \leq \lt[\prod_{i=1}^n(d_i-1) - \lt(1-q^{-\nicefrac12}\rt)\cdot M(d_1,\dots, d_n)\rt]q^{\nicefrac{(n-2)}{2}}.$$
\end{Theorem}
\begin{proof}
  By \JacobiCompletethm[(iii-iv)] if $\alpha\neq 0$ then $|J_{\alpha}(\lm_1^{t_1},\dots, \lm_n^{t_n})|=q^{\nicefrac{(n-1)}{2}}$. So we have $$\lt|Z\lt(a_1X_1^{k_1}+\cdots a_nX_n^{k_n}=\alpha\rt)-q^{n-1}\rt|\leq \lt[\prod_{i=1}^n(d_i-1)-|T|\rt]q^{\nicefrac{(n-1)}{2}}+|T|\cdot q^{\nicefrac{(q-2)}{2}}=\lt[\prod_{i=1}^n (d_i-1)-\lt(1-q^{-\nicefrac12}\rt)|T|\rt]q^{\nicefrac{(n-1)}{2}}$$Since $|T|=M(d_1,\dots, d_n)$ we have the theorem. 
\end{proof}

\subsection{A General Formula of \texorpdfstring{$M(d_1,\dots, d_n)$}{M(d1,...,dn)}}
We follow the notation from the previous subsection. Now suppose one of the $d_i$ is relatively prime to all other $d_j$'s where $j\neq i$ i.e. $\gcd(d_i,d_j)=1$ for all $j\in[n]$, $j\neq i$. If for some $(t_1,\dots, t_n)\in\bigtimes_{i=1}^n[d_i-1]$ we have $\sum_{i=1}^n \frac{t_i}{d_i}=m\in\bbZ$ then we have $$t_1\cdot (d_2\cdots d_n)+k\cdot d_1=m\cdot (d_1,c\dots d_n)$$Thus $t_1\cdot (d_2\cdots d_n)\equiv 0\bmod{d_1}$. Since $d_1$ is relatively prime with all other $d_j$ for all $j\in\{2,\dots, n\}$ we have $d_1\mid t_1$. Hence contradiction \ctr Hence such tuple $(t_1,\dots,t_n)$ does not exist. Therefore $M(d_1,\dots, d_n)=0$. 

Now we already have a trivial upper bound on $M(d_1,\dots, d_n)\leq \prod_{i=1}^n (d_i-1)$. We will give a general formula for $M(d_1,\dots, d_n)$. Suppose $D=\lcm(d_1,\dots, d_n)$. Then for any tuple $(t_1,\dots, t_n)\in \bigtimes_{i=1}^n[d_i-1]$ observe that  \begin{equation}\label{eq:integrality-exponential-sum}
  \frac1D\sum_{h=0}^{D-1}\exp\lt[2\pi ih\lt(\frac{t_1}{d_1}+\cdots+\frac{t_n}{d_n}\rt)\rt]=\begin{cases}
    1 & \text{if $\frac{t_1}{d_1}+\cdots+\frac{t_n}{d_n}\in\bbZ$}\\ 
    0 & \text{otherwise}
  \end{cases}
\end{equation}
Then using this expression we have \begin{align*}
  M(d_1,\dots, d_n) & = \sum_{t_1=1}^{d_1-1}\cdots \sum_{t_n=1}^{d_n-1}\frac1D\sum_{h=0}^{D-1}\exp\lt[2\pi ih\sum_{i=1}^n\frac{t_i}{d_i}\rt]\\ 
    & = \frac1D\sum_{h=0}^{D-1} \lt(\;\sum_{t_1=1}^{d_1-1}\exp\lt[t_1\cdot \frac{h}{d_1}\rt]\;\rt)\cdots \lt(\;\sum_{t_n=1}^{d_n-1}\exp\lt[t_n\cdot \frac{h}{d_n}\rt]\;\rt)\\ 
    & = \frac1D\sum_{h=0}^{D-1}\prod_{i=1}^n \lt(\;\sum_{t_i=1}^{d_i-1}\exp\lt[t_i\cdot \frac{h}{d_i}\rt]\;\rt)\\ 
    & = \frac1D\sum_{h=0}^{D-1}\prod_{i=1}^n \lt(\;\sum_{t_i=0}^{d_i-1}\exp\lt[t_i\cdot \frac{h}{d_i}\rt]-1\;\rt)\\ 
    & = \frac1D\sum_{h=0}^{D-1}\lt[ (-1)^n+\sum_{k=1}^n(-1)^{n-k}\sum_{1\leq i_1<\cdots <i_k\leq n}\lt( \sum_{t_{i_1}=0}^{d_{i_1}-1}\exp\lt[ t_{i_1}\frac{h}{d_{i_1}} \rt] \rt)\cdots   \lt( \sum_{t_{i_n}=0}^{d_{i_n}-1}\exp\lt[ t_{i_n}\frac{h}{d_{i_n}} \rt] \rt) \rt]\\ 
    & = (-1)^n+\frac1D\sum_{h=0}^{D-1} \sum_{k=1}^n(-1)^{n-k}\sum_{1\leq i_1<\cdots <i_k\leq n}\prod_{j=1}^k\lt(\; \sum_{t_{i_j}=0}^{d_{i_j}-1}\exp\lt[ t_{i_j}\frac{h}{d_{i_j}} \rt]\; \rt)
\end{align*}
Now for any $d\in\bbN$ we have $$\sum_{j=0}^d\exp\lt[j\frac{h}d\rt]=\begin{cases}
  d & \text{if $h\equiv 0\bmod d$}\\ 
  0 & \text{otherwise}
\end{cases}$$So in the product term at the end of the last expression survives if $d_{i_j}\mid h$ for all $j\in[k]$ i.e. $\lcm(d_{i_1},\cdots, d_{i_k})\mid h$. So in the product at the end of last expression we obtain 
$$\prod_{j=1}^k\lt(\; \sum_{t_{i_j}=0}^{d_{i_j}-1}\exp\lt[ t_{i_j}\frac{h}{d_{i_j}} \rt]\rt)=\begin{cases}
  d_{i_1}\cdots d_{i_k}& \text{if $d_{i_j}\mid h$ for all $j\in[k]$}\\ 
  0 & \text{otherwise}
\end{cases}$$Hence we have \begin{align*}
  M(d_1,\dots, d_n) & = (-1)^n+\frac1D\sum_{h=0}^{D-1} \sum_{k=1}^n(-1)^{n-k}\sum_{\substack{1\leq i_1<\cdots <i_k\leq n\\[1pt] \lcm\lt(d_{i_1},\dots,d_{i_k}\rt)\mid h}}\lt(d_{i_1}\cdots d_{i_k}\rt)\\ 
  & = (-1)^n+\sum_{h=0}^{D-1} \sum_{k=1}^n(-1)^{n-k}\sum_{1\leq i_1<\cdots <i_k\leq n}\lt(d_{i_1}\cdots d_{i_k}\rt)\frac1D\sum_{\substack{h=0\\[1pt] \lcm\lt(d_{i_1},\dots, d_{i_n}\rt)\mid h}}^{D-1}1\\ 
  & = (-1)^n+\sum_{h=0}^{D-1} \sum_{k=1}^n(-1)^{n-k}\sum_{1\leq i_1<\cdots <i_k\leq n}\frac{d_{i_1}\cdots d_{i_k}}{\lcm\lt(d_{i_1},\dots, d_{i_n}\rt)}
\end{align*}
So we have the following theorem summarizing the formula.
\begin{Theorem}{Closed Form for $M(d_1,\dots,d_n)$}{thm:M-formula}
  Let $d_1,\dots,d_n\in\bbN$ such that $d_i\mid q-1$ for all $i\in[n]$ and set $D=\lcm(d_1,\dots,d_n)$. Then
  $$M(d_1,\dots,d_n)=(-1)^n+\sum_{k=1}^n(-1)^{n-k}\sum_{1\leq i_1<\cdots<i_k\leq n}\frac{d_{i_1}\cdots d_{i_k}}{\lcm\lt(d_{i_1},\dots,d_{i_k}\rt)}.$$
\end{Theorem}
